% SPDX-License-Identifier: MIT
% Copyright (c) 2026 Evgeny Baulin
% Optimization Methods, Seminar 03: theory handout (Russian).
% Generated by checks/03. Gradient Descent and Automatic Differentiation.ipynb from web/03/js/content-*.js and content/students.json. Do not edit by hand.
% Compile with: latexmk -xelatex -interaction=nonstopmode -halt-on-error Theory_ru.tex
\documentclass[11pt,a4paper]{article}

\usepackage{fontspec}
\setmainfont{cmunrm.otf}[BoldFont=cmunbx.otf, ItalicFont=cmunti.otf, BoldItalicFont=cmunbi.otf]
\setsansfont{cmunss.otf}[BoldFont=cmunsx.otf, ItalicFont=cmunsi.otf, BoldItalicFont=cmunso.otf]
\setmonofont{cmuntt.otf}[BoldFont=cmuntb.otf, ItalicFont=cmunit.otf, BoldItalicFont=cmuntx.otf]
\usepackage{polyglossia}
\setdefaultlanguage{russian}
\setotherlanguage{english}
\newfontfamily\cyrillicfont{cmunrm.otf}[Ligatures=TeX, BoldFont=cmunbx.otf, ItalicFont=cmunti.otf, BoldItalicFont=cmunbi.otf]
\newfontfamily\cyrillicfontsf{cmunss.otf}[Ligatures=TeX, BoldFont=cmunsx.otf, ItalicFont=cmunsi.otf, BoldItalicFont=cmunso.otf]
\newfontfamily\cyrillicfonttt{cmuntt.otf}[BoldFont=cmuntb.otf, ItalicFont=cmunit.otf, BoldItalicFont=cmuntx.otf]

\usepackage[a4paper,left=2.3cm,right=2.3cm,top=2.4cm,bottom=2.4cm,headheight=14pt]{geometry}
\usepackage{amsmath,amssymb}
\usepackage{graphicx}
\usepackage{array,booktabs,longtable}
\usepackage{enumitem}
\usepackage{xcolor}
\usepackage[most]{tcolorbox}
\usepackage{titlesec}
\usepackage{fancyhdr}
\usepackage{needspace}
\usepackage[expansion=false]{microtype}
\usepackage[colorlinks=true,linkcolor=hseblue,urlcolor=hseblue,citecolor=hseblue]{hyperref}

\definecolor{hseblue}{RGB}{16,45,105}
\definecolor{hseaccent}{RGB}{198,40,40}
\definecolor{hsegray}{RGB}{108,114,126}
\definecolor{hselight}{RGB}{240,243,248}
\definecolor{hsewarm}{RGB}{253,244,232}

\hypersetup{pdftitle={Методы оптимизации. Семинар 3. Конспект теории}, pdfauthor={Евгений Баулин}}

\titleformat{\section}{\color{hseblue}\Large\bfseries}{\thesection}{0.6em}{}
\titleformat{\subsection}{\color{hseblue}\large\bfseries}{\thesubsection}{0.6em}{}
\titlespacing*{\section}{0pt}{16pt plus 4pt minus 2pt}{8pt plus 2pt}
\titlespacing*{\subsection}{0pt}{12pt plus 3pt minus 2pt}{6pt plus 2pt}
% a heading keeps a few lines under it on its page: a breakable box right after it would otherwise leave it alone
\pretocmd{\section}{\Needspace{7\baselineskip}}{}{\errmessage{cannot patch section}}
\pretocmd{\subsection}{\Needspace{5\baselineskip}}{}{\errmessage{cannot patch subsection}}

\tcbset{handoutbox/.style={enhanced, breakable, sharp corners, boxrule=0pt, left=7pt, right=7pt, top=5pt, bottom=5pt, before skip=8pt plus 2pt, after skip=8pt plus 2pt}}
\newtcolorbox{examplebox}{handoutbox, colback=hselight, colframe=hselight, before upper={\textcolor{hseblue}{\bfseries Пример.}\ }}
\newtcolorbox{whybox}{handoutbox, colback=white, colframe=hseblue, leftrule=2.5pt, top=3pt, bottom=3pt, before upper={\textcolor{hseaccent}{\bfseries Почему.}\ }}
\newtcolorbox{answerbox}{handoutbox, colback=hsewarm, colframe=hsewarm, before upper={\textcolor{hseaccent}{\bfseries Ответ.}\ }}
\newtcolorbox{conventionbox}{handoutbox, colback=hselight, colframe=hseblue, leftrule=2.5pt, before upper={\textcolor{hseblue}{\bfseries Обозначения.}\par\smallskip}}

\newcommand{\handoutfigure}[3]{\par\medskip\noindent\begin{minipage}{\linewidth}\centering
  \includegraphics[width=#1\linewidth]{figures/#2.pdf}\par\smallskip
  \begin{minipage}{0.94\linewidth}\small\textcolor{hseblue}{\bfseries Рисунок.}\ #3\end{minipage}\end{minipage}\par\medskip}
\newcommand{\problemhead}[3]{\par\bigskip\noindent{\color{hseblue}\large\bfseries #1.\ #2}\par\nopagebreak
  \noindent{\small\itshape\color{hsegray}#3}\par\nopagebreak\smallskip}
\newtcolorbox{definitionbox}[1]{handoutbox, colback=white, colframe=hseblue, leftrule=2.5pt, before upper={\textcolor{hseblue}{\bfseries #1}\ }}
\newtcolorbox{theorembox}[1]{handoutbox, colback=hselight, colframe=hseblue, leftrule=2.5pt, before upper={\textcolor{hseblue}{\bfseries #1}\ }}
\newtcolorbox{proofbox}{handoutbox, colback=white, colframe=hsegray, leftrule=1.2pt, top=3pt, bottom=3pt, before upper={\textcolor{hsegray}{\bfseries\itshape Доказательство.}\ }}
\newtcolorbox{notebox}[1]{handoutbox, colback=hselight, colframe=hselight, before upper={\textcolor{hseblue}{\bfseries #1}\par\smallskip}}
\newtcolorbox{trapbox}{handoutbox, colback=hsewarm, colframe=hseaccent, leftrule=2.5pt, before upper={\textcolor{hseaccent}{\bfseries Ловушка.}\ }}
\newtcolorbox{selfcheckbox}{handoutbox, colback=white, colframe=hseaccent, leftrule=2.5pt, top=3pt, bottom=3pt, before upper={\textcolor{hseaccent}{\bfseries Проверьте себя.}\ }}
% a display wider than the line is set as a centred paragraph that may break after a relation
\newcommand{\handoutdisplay}[1]{\sbox0{$\displaystyle #1$}\ifdim\wd0>\linewidth
  \par{\centering\binoppenalty=10000\relpenalty=0\relax$\displaystyle #1$\par}\else\[#1\]\fi}
\newcommand{\selfcheckanswer}{\par\smallskip\textcolor{hsegray}{\bfseries Ответ.}\ }

\setlist[itemize]{itemsep=3pt,topsep=4pt}
\setlist[enumerate]{itemsep=3pt,topsep=4pt}
\setlength{\parindent}{0pt}
\setlength{\parskip}{5pt plus 1pt}
\setlength{\emergencystretch}{2em}
\setcounter{tocdepth}{2}

\pagestyle{fancy}
\fancyhf{}
\fancyhead[L]{\small\color{hsegray}Методы оптимизации}
\fancyhead[R]{\small\color{hsegray}Семинар 3}
\fancyfoot[L]{\small\color{hsegray}Евгений Баулин}
\fancyfoot[C]{\small\thepage}
\fancyfoot[R]{\small\color{hsegray}Лицензия MIT}
\renewcommand{\headrulewidth}{0pt}
\renewcommand{\footrulewidth}{0pt}
\fancypagestyle{plain}{\pagestyle{fancy}}

\begin{document}

\begin{center}
  {\large\color{hsegray}Методы оптимизации\par}
  \vspace{4pt}
  {\large\bfseries\color{hseaccent}Семинар 3\par}
  \vspace{8pt}
  {\LARGE\bfseries\color{hseblue}Градиентный спуск и автоматическое дифференцирование\par}
  \vspace{8pt}
  {\Large Конспект теории\par}
  \vspace{12pt}
  НИУ ВШЭ, факультет компьютерных наук\par
  2026–2027 учебный год\par
  \vspace{8pt}
  Лектор курса: Андрей Игнатов\par
  Евгений Баулин\par
  Telegram: \href{https://t.me/tarakan_tuc}{@tarakan\_tuc}\par
\end{center}

\noindent{\color{hseblue}\rule{\linewidth}{0.6pt}}

\tableofcontents

\clearpage


Этот конспект собирает теорию третьего семинара: определения для коллоквиума, неравенство Йенсена, градиентный спуск с постоянным шагом и с бэктрекингом, скорости сходимости и автоматическое дифференцирование в обратном режиме. Каждое утверждение дано вместе с условиями, а там, где его могут спросить на коллоквиуме, — с полным доказательством. Задачи семинара разобраны целиком.

\begin{conventionbox}
Векторы — столбцы, $x^\top y$ — скалярное произведение, $\|\cdot\|$ — евклидова норма, если нижний индекс не указывает иное. Функция $f$ задана на $\mathbb R^n$, $f^* = \inf_x f(x)$ — оптимальное значение, $x^*$ — точка минимума, если она существует. Итерации — $x_0, x_1, \dots$, шаг — $\alpha$ или $\alpha_k$, направление — $d$ или $d_k$, а $x^+$ внутри доказательства — следующая итерация. Константы сильной выпуклости и гладкости — $\mu$ и $L$, а $\kappa = L/\mu$. Для симметричных матриц запись $A \preceq B$ означает, что матрица $B - A$ положительно полуопределена. В автоматическом дифференцировании $\dot v$ — касательная переменная, а $\bar v = \partial z/\partial v$ — сопряжённая переменная.
\end{conventionbox}

% =====================================================================
\section{Определения}\label{sec:definitions}

На коллоквиуме сначала спрашивают определения, и определение засчитывается только вместе с условиями. Каждый блок ниже — формулировка, которую нужно воспроизвести; короткие абзацы между блоками — следствия, а не часть определений.

\begin{definitionbox}{Определение: Локальный и глобальный минимум.}
Пусть $f: \mathbb R^n \to \mathbb R$. Точка $x^*$ — \textbf{точка глобального минимума} $f$, если $f(x^*) \le f(x)$ для всех $x \in \mathbb R^n$. Это \textbf{точка локального минимума}, если существует $\varepsilon > 0$, при котором $f(x^*) \le f(x)$ для всех $x$ с $\|x - x^*\| < \varepsilon$. Точка минимума \textbf{строгая}, если неравенство строгое для всех $x \ne x^*$ из соответствующего множества. Тогда значение $f(x^*)$ называют глобальным или локальным \textbf{минимумом}, а $f^* = \inf_x f(x)$ — оптимальным значением.
\end{definitionbox}

Оптимальное значение может равняться $-\infty$, а может быть конечным, но не достигаться: у $f(x) = e^x$ на $\mathbb R$ значение $f^* = 0$, а точки минимума нет.

\begin{definitionbox}{Определение: Стационарная точка.}
Пусть $f: \mathbb R^n \to \mathbb R$ дифференцируема в точке $x$. Точка $x$ называется \textbf{стационарной}, если $\nabla f(x) = 0$.
\end{definitionbox}

Точка локального минимума дифференцируемой функции стационарна: если $\nabla f(x^*) \ne 0$, то $d = -\nabla f(x^*)$ — направление спуска (определение ниже), и вдоль него $f$ убывает. Обратное неверно: стационарная точка может оказаться точкой локального максимума или седловой точкой. У выпуклой дифференцируемой $f$ любая стационарная точка — точка глобального минимума: достаточно подставить $\nabla f(x^*) = 0$ в неравенство для касательной из следующего раздела.

\begin{definitionbox}{Определение: Выпуклое множество.}
Множество $C \subseteq \mathbb R^n$ \textbf{выпукло}, если $(1-t)x + ty \in C$ для всех $x, y \in C$ и всех $t \in [0, 1]$, то есть если вместе с любыми двумя своими точками оно содержит отрезок между ними.
\end{definitionbox}

\begin{definitionbox}{Определение: Выпуклая и строго выпуклая функция.}
Пусть $C \subseteq \mathbb R^n$ — выпуклое множество, обычно $C = \mathbb R^n$. Функция $f: C \to \mathbb R$ \textbf{выпукла}, если для неё выполнено неравенство Йенсена
\handoutdisplay{
f\big((1-t)x + ty\big) \le (1-t)f(x) + t f(y) \quad \text{для всех } x, y \in C,\ t \in [0, 1].
}
 Она \textbf{строго выпукла}, если неравенство строгое при $x \ne y$ и $t \in (0, 1)$, и \textbf{вогнута}, если функция $-f$ выпукла.
\end{definitionbox}

\begin{definitionbox}{Определение: $\mu$-сильная выпуклость.}
Пусть $\mu > 0$. Функцию $f: \mathbb R^n \to \mathbb R$ называют $\mu$-\textbf{сильно выпуклой}, если функция $f(x) - \frac{\mu}{2}\|x\|^2$ выпукла. Равносильное определение — неравенство Йенсена с квадратичным запасом:
\handoutdisplay{
f\big((1-t)x + ty\big) \le (1-t)f(x) + t f(y) - \frac{\mu}{2}\,t(1-t)\|x - y\|^2 \quad \text{для всех } x, y,\ t \in [0, 1].
}
 Если $f$ дифференцируема, это равносильно условию первого порядка
\handoutdisplay{
f(y) \ge f(x) + \nabla f(x)^\top (y - x) + \frac{\mu}{2}\|y - x\|^2 \quad \text{для всех } x, y,
}
 а если $f$ дважды дифференцируема — условию $\nabla^2 f(x) \succeq \mu I$ для всех $x$.
\end{definitionbox}

Эти формулировки равносильны, потому что $(1-t)\|x\|^2 + t\|y\|^2 - \|(1-t)x + ty\|^2 = t(1-t)\|x - y\|^2$ (тождество из следующего раздела), а условие первого порядка — это неравенство для касательной, записанное для $f - \frac{\mu}{2}\|\cdot\|^2$. Сильно выпуклая функция строго выпукла и имеет ровно одну точку минимума. Строгая выпуклость слабее: у $f(x) = \log(1 + e^x)$ вторая производная $f''(x) = \sigma(x)\big(1 - \sigma(x)\big) > 0$, где $\sigma(x) = 1/(1 + e^{-x})$, но $f''(x) \to 0$ при $|x| \to \infty$, так что никакое $\mu > 0$ не подходит, а точки минимума у этой $f$ нет вовсе.

\begin{definitionbox}{Определение: $L$-гладкость.}
Пусть $L > 0$. Дифференцируемую функцию $f: \mathbb R^n \to \mathbb R$ называют $L$-\textbf{гладкой}, если её градиент $L$-липшицев:
\handoutdisplay{
\|\nabla f(x) - \nabla f(y)\| \le L\|x - y\| \quad \text{для всех } x, y.
}
 Если $f$ дважды непрерывно дифференцируема, это равносильно условию $-LI \preceq \nabla^2 f(x) \preceq LI$ для всех $x$. Для выпуклой $f$ левое неравенство выполнено автоматически, и условие принимает вид $\nabla^2 f(x) \preceq LI$.
\end{definitionbox}

В лекции одно и то же свойство называется двумя способами: $L$-гладкость и $L$-липшицев градиент. Доказательства сходимости используют его в форме леммы о спуске из раздела о градиентном спуске: $f(y) \le f(x) + \nabla f(x)^\top (y - x) + \frac{L}{2}\|y - x\|^2$.

\begin{definitionbox}{Определение: Число обусловленности.}
Пусть $f$ — $\mu$-сильно выпуклая и $L$-гладкая функция, причём $\mu$ — наибольшая, а $L$ — наименьшая из констант, при которых это верно. \textbf{Число обусловленности} $f$ — это $\kappa = L/\mu$, и $\kappa \ge 1$. У квадратичной функции $f(x) = \tfrac{1}{2} x^\top Q x - b^\top x$ с симметричной матрицей $Q \succ 0$ наилучшие константы — $\mu = \lambda_{\min}(Q)$ и $L = \lambda_{\max}(Q)$, поэтому $\kappa = \lambda_{\max}(Q)/\lambda_{\min}(Q)$.
\end{definitionbox}

Множества уровня такой квадратичной функции — эллипсоиды, и отношение их наибольшей полуоси к наименьшей равно $\sqrt{\kappa}$. Большое $\kappa$ означает длинную узкую долину.

\begin{definitionbox}{Определение: Направление спуска.}
Пусть $f: \mathbb R^n \to \mathbb R$ дифференцируема в точке $x$. Вектор $d \in \mathbb R^n$ называется \textbf{направлением спуска} в точке $x$, если $\nabla f(x)^\top d < 0$. Тогда $f(x + \alpha d) < f(x)$ при всех достаточно малых $\alpha > 0$, поскольку $f(x + \alpha d) = f(x) + \alpha\,\nabla f(x)^\top d + o(\alpha)$.
\end{definitionbox}

Если $\nabla f(x) \ne 0$, то вектор $d = -\nabla f(x)$ — направление спуска: $\nabla f(x)^\top d = -\|\nabla f(x)\|^2 < 0$. Это условие достаточно для локального убывания, но не необходимо: в точке максимума $0$ функции $f(x) = -x^2$ любое $d \ne 0$ уменьшает $f$, хотя $f'(0)\,d = 0$.

\begin{definitionbox}{Определение: Условие Поляка–Лоясевича.}
Пусть $f: \mathbb R^n \to \mathbb R$ дифференцируема, $f^* = \inf_x f(x) > -\infty$ и $\mu > 0$. Функция $f$ удовлетворяет \textbf{условию Поляка–Лоясевича (PL)} с константой $\mu$, если
\handoutdisplay{
\frac{1}{2}\|\nabla f(x)\|^2 \ge \mu\big(f(x) - f^*\big) \quad \text{для всех } x.
}

\end{definitionbox}

При условии PL любая стационарная точка — точка глобального минимума: из $\nabla f(x) = 0$ следует $f(x) = f^*$. Дифференцируемая $\mu$-сильно выпуклая функция удовлетворяет PL с тем же $\mu$: минимизируя правую часть условия первого порядка по $y$ (минимум достигается при $y = x - \nabla f(x)/\mu$), получаем $f^* \ge f(x) - \frac{1}{2\mu}\|\nabla f(x)\|^2$. Обратное неверно, и PL не требует даже выпуклости: $f(x) = x^2 + 3\sin^2 x$ удовлетворяет PL с некоторым $\mu > 0$, хотя $f''(\pi/2) = 2 + 6\cos\pi < 0$.

% =====================================================================
\section{Неравенство Йенсена}\label{sec:jensen}

Неравенство Йенсена — это определение выпуклости. Обычные условия теорем о сходимости — то же неравенство с квадратичным запасом: снизу для сильной выпуклости, сверху для гладкости.

\subsection{Как читать неравенство}\label{sub:jensen-reading}

Выпуклость $f$ на выпуклом множестве — это неравенство Йенсена
\handoutdisplay{
f\big((1-t)x + ty\big) \le (1-t)f(x) + t f(y), \qquad t \in [0, 1].
}
 Когда $t$ меняется от $0$ до $1$, точка $(1-t)x + ty$ пробегает отрезок от $x$ до $y$, так что левая часть — это график $f$ над отрезком. В правой части те же веса применены к значениям функции: это хорда от $(x, f(x))$ до $(y, f(y))$. Выпуклость означает ровно то, что хорда нигде не опускается ниже графика. Производные здесь не участвуют: $f(x) = |x|$ выпукла в силу неравенства треугольника $|(1-t)x + ty| \le (1-t)|x| + t|y|$, хотя в точке $0$ производной у неё нет.

Ограничение $t \in [0, 1]$ существенно. При других $t$ точка $(1-t)x + ty$ по-прежнему лежит на прямой, проходящей через $x$ и $y$, но уже не между ними, и неравенство нарушается даже для $f(x) = x^2$. При $x = 0$, $y = 1$ и $t = 2$ получаем точку $2$; левая часть равна $f(2) = 4$, а правая — $(1-t)f(x) + tf(y) = 2$.

\begin{theorembox}{Утверждение: Тождество для $x^2$.}
Для всех вещественных $x$, $y$ и $t$
\handoutdisplay{
(1-t)x^2 + ty^2 - \big((1-t)x + ty\big)^2 = t(1-t)(x - y)^2.
}

\end{theorembox}

\begin{proofbox}
Раскроем квадрат $\big((1-t)x + ty\big)^2 = (1-t)^2x^2 + 2t(1-t)xy + t^2y^2$ и воспользуемся равенствами $(1-t) - (1-t)^2 = t(1-t)$ и $t - t^2 = t(1-t)$: разность равна $t(1-t)(x^2 - 2xy + y^2) = t(1-t)(x - y)^2$.
\end{proofbox}

При $t \in [0, 1]$ правая часть неотрицательна, поэтому $x^2$ выпукла, а в нуль она обращается только при $t \in \{0, 1\}$ или $x = y$, поэтому $x^2$ строго выпукла. При $t \notin [0, 1]$ множитель $t(1-t)$ отрицателен, и неравенство нарушается для любых $x \ne y$; в примере выше $t(1-t)(x - y)^2 = 2 - 4 < 0$. То же вычисление со скалярными произведениями даёт для векторов
\handoutdisplay{
(1-t)\|x\|^2 + t\|y\|^2 - \|(1-t)x + ty\|^2 = t(1-t)\|x - y\|^2,
}
 а для $f(x) = \tfrac{1}{2} x^\top Q x$ с симметричной $Q$ зазор между хордой и графиком равен $\tfrac{1}{2}\,t(1-t)(x - y)^\top Q (x - y)$.

\subsection{Три равносильных критерия}\label{sub:jensen-characterizations}

Определение выпуклости — неравенство для хорды. Ему равносильны ещё два критерия, но каждый из них требует гладкости, а определение обходится без неё.

\begin{center}\small
\begin{tabular}{@{}>{\raggedright\arraybackslash\hspace{0pt}}p{0.24\linewidth}>{\raggedright\arraybackslash\hspace{0pt}}p{0.25\linewidth}>{\raggedright\arraybackslash\hspace{0pt}}p{0.4\linewidth}@{}}
\toprule
Критерий & Что требуется & Смысл \\
\midrule
Хорда (Йенсен): $f\big((1-t)x + ty\big) \le (1-t)f(x) + tf(y)$ & выпуклая область определения, производные не нужны & график лежит не выше любой хорды \\
Касательная: $f(y) \ge f(x) + \nabla f(x)^\top (y - x)$ & $f$ дифференцируема на открытом выпуклом множестве & график лежит не ниже любой касательной плоскости; линейная модель — глобальная оценка снизу \\
Кривизна: $\nabla^2 f(x) \succeq 0$ & $f$ дважды дифференцируема на открытом выпуклом множестве & $d^\top \nabla^2 f(x)\, d \ge 0$ для любого $d$: вдоль любой прямой график изгибается вверх \\
\bottomrule
\end{tabular}
\end{center}

\begin{theorembox}{Теорема: Критерии выпуклости.}
Пусть $C \subseteq \mathbb R^n$ — открытое выпуклое множество. Если $f$ дифференцируема на $C$, то она выпукла тогда и только тогда, когда $f(y) \ge f(x) + \nabla f(x)^\top (y - x)$ для всех $x, y \in C$. Если $f$ дважды дифференцируема на $C$, то она выпукла тогда и только тогда, когда $\nabla^2 f(x) \succeq 0$ для всех $x \in C$.
\end{theorembox}

\begin{proofbox}
От хорды к касательной: неравенство Йенсена при $t \in (0, 1]$ даёт $f\big(x + t(y - x)\big) - f(x) \le t\big(f(y) - f(x)\big)$; делим на $t$ и, переходя к пределу при $t \to 0$, получаем $\nabla f(x)^\top (y - x) \le f(y) - f(x)$. От касательной к хорде: для $z = (1-t)x + ty$ неравенство для касательной в точке $z$ даёт $f(x) \ge f(z) + \nabla f(z)^\top (x - z)$ и $f(y) \ge f(z) + \nabla f(z)^\top (y - z)$; умножаем их на $1-t$ и $t$ соответственно и складываем; слагаемые с градиентом сокращаются, так как $(1-t)(x - z) + t(y - z) = 0$. Кривизна: $f$ выпукла тогда и только тогда, когда выпукло каждое её сужение $\varphi(s) = f(x + sd)$ на прямую; дважды дифференцируемая функция одной переменной выпукла тогда и только тогда, когда $\varphi'' \ge 0$, а $\varphi''(s) = d^\top \nabla^2 f(x + sd)\, d$.
\end{proofbox}

Для строгой выпуклости верна только одна импликация: из $\nabla^2 f(x) \succ 0$ при всех $x$ следует строгая выпуклость, но $f(x) = |x|^3$ строго выпукла, хотя $f''(x) = 6|x|$ обращается в нуль в точке $0$.

\subsection{Конечная форма и форма для математического ожидания}\label{sub:jensen-forms}

\begin{theorembox}{Теорема: Неравенство Йенсена для нескольких точек.}
Пусть $f$ выпукла на выпуклом множестве $C$, $x_1, \dots, x_N \in C$ и $t_1, \dots, t_N \ge 0$, причём $\sum_i t_i = 1$. Тогда
\handoutdisplay{
f\Big(\sum_{i=1}^N t_i x_i\Big) \le \sum_{i=1}^N t_i f(x_i).
}

\end{theorembox}

\begin{proofbox}
Индукция по $N$. При $N = 1$ обе части равны $f(x_1)$, а случай $N = 2$ — это определение. Пусть $N \ge 3$ и утверждение верно для $N - 1$ точки. Если $t_N = 1$, все остальные веса равны нулю и доказывать нечего. Иначе
\handoutdisplay{
\sum_{i=1}^N t_i x_i = (1 - t_N)\,y + t_N x_N, \qquad y = \sum_{i=1}^{N-1} \frac{t_i}{1 - t_N}\,x_i,
}
 и $y$ — выпуклая комбинация $N - 1$ точки из $C$: её веса неотрицательны и в сумме дают $1$, поэтому $y \in C$. Определение, а затем предположение индукции дают
\handoutdisplay{
f\Big(\sum_{i=1}^N t_i x_i\Big) \le (1 - t_N) f(y) + t_N f(x_N) \le \sum_{i=1}^{N-1} t_i f(x_i) + t_N f(x_N).
}

\end{proofbox}

\begin{theorembox}{Теорема: Неравенство Йенсена для математического ожидания.}
Пусть $f: \mathbb R^n \to \mathbb R$ выпукла, а $X$ — случайный вектор в $\mathbb R^n$, для которого $\mathbb E\|X\| < \infty$ и $\mathbb E|f(X)| < \infty$. Тогда
\handoutdisplay{
f(\mathbb E X) \le \mathbb E f(X).
}

\end{theorembox}

\begin{proofbox}
Положим $m = \mathbb E X$. Если $f$ дифференцируема, неравенство для касательной в точке $m$ имеет вид $f(X) \ge f(m) + \nabla f(m)^\top (X - m)$; берём математическое ожидание обеих частей и пользуемся тем, что $\mathbb E(X - m) = 0$. Без дифференцируемости то же рассуждение проходит с субградиентом в точке $m$ — вектором $g$, для которого $f(y) \ge f(m) + g^\top (y - m)$ при всех $y$; у выпуклой $f: \mathbb R^n \to \mathbb R$ он существует в каждой точке. Случайный вектор, принимающий значение $x_i$ с вероятностью $t_i$, возвращает нас к конечной форме.
\end{proofbox}

\subsection{Вогнутые функции и неравенство о средних}\label{sub:jensen-concave}

\begin{trapbox}
Для вогнутой $g$ знак каждого неравенства выше меняется на противоположный: $g\big(\sum_i t_i x_i\big) \ge \sum_i t_i g(x_i)$ и $g(\mathbb E X) \ge \mathbb E g(X)$. Логарифм вогнут, поэтому для положительной случайной величины $X$ выполнено $\mathbb E \log X \le \log \mathbb E X$, а не наоборот.
\end{trapbox}

\begin{theorembox}{Утверждение: Неравенство о средних для двух чисел.}
Для $a, b > 0$
\handoutdisplay{
\sqrt{ab} \le \frac{a + b}{2},
}
 причём равенство достигается тогда и только тогда, когда $a = b$.
\end{theorembox}

\begin{proofbox}
Функция $-\log$ строго выпукла на выпуклом множестве $(0, \infty)$, так как $(-\log x)'' = 1/x^2 > 0$. Неравенство Йенсена при $t = \tfrac{1}{2}$ даёт
\handoutdisplay{
-\log\frac{a + b}{2} \le \frac{1}{2}\big(-\log a - \log b\big) = -\log\sqrt{ab},
}
 причём неравенство строгое при $a \ne b$. Остаётся умножить на $-1$ и применить возрастающую функцию $\exp$.
\end{proofbox}

\begin{examplebox}
$a = 1$, $b = 9$: $\sqrt{ab} = 3 \le 5 = \frac{a + b}{2}$. Конечная форма неравенства Йенсена тем же рассуждением даёт взвешенный вариант $\prod_i a_i^{t_i} \le \sum_i t_i a_i$ для $a_i > 0$, $t_i \ge 0$, $\sum_i t_i = 1$.
\end{examplebox}

\subsection{Контрпример}\label{sub:jensen-counterexample}

\begin{examplebox}
Возьмём $f(x) = x^3$ на $\mathbb R$ и $x = -2$, $y = 1$, $t = 0{,}5$. Тогда $(1-t)x + ty = -0{,}5$, левая часть равна $f(-0{,}5) = -0{,}125$, а правая — $(1-t)x^3 + ty^3 = -3{,}5$. Так как $-0{,}125 > -3{,}5$, неравенство нарушено, и $x^3$ не выпукла на $\mathbb R$.
\end{examplebox}

Одной пары точек и одного $t$ достаточно, чтобы опровергнуть выпуклость; никакое число удачных проверок её не доказывает. На $[0, \infty)$ функция $x^3$ выпукла, так как там $f''(x) = 6x \ge 0$, поэтому случайная проверка, которая берёт 10000 пар из $[0, 2]$, находит 0 нарушений и ничего не говорит об отрицательных аргументах. Выпуклость доказывают по определению или с помощью одного из критериев выше.

\subsection{Квадратичный запас с обеих сторон}\label{sub:jensen-margins}

На языке хорды $\mu$ и $L$ ограничивают зазор между хордой и графиком снизу и сверху. Если $f$ — $\mu$-сильно выпуклая и $L$-гладкая функция, то для всех $x, y$ и $t \in [0, 1]$
\handoutdisplay{
\frac{\mu}{2}\,t(1-t)\|x - y\|^2 \le (1-t)f(x) + tf(y) - f\big((1-t)x + ty\big) \le \frac{L}{2}\,t(1-t)\|x - y\|^2,
}
 а в форме первого порядка
\handoutdisplay{
\frac{\mu}{2}\|y - x\|^2 \le f(y) - f(x) - \nabla f(x)^\top (y - x) \le \frac{L}{2}\|y - x\|^2.
}
 Левые неравенства — определение сильной выпуклости. Правое неравенство в форме первого порядка — лемма о спуске, а его вариант для хорды получается так: записываем лемму о спуске в точке $z = (1-t)x + ty$ для точек $x$ и $y$ и складываем с весами $1-t$ и $t$. Для выпуклой дифференцируемой $f$ каждое из правых неравенств на самом деле равносильно $L$-гладкости.

Для $f(x) = \tfrac{1}{2} x^\top Q x$ зазор между хордой и графиком равен $\tfrac{1}{2}\,t(1-t)(x - y)^\top Q (x - y)$. Он лежит между двумя оценками с $\mu = \lambda_{\min}(Q)$ и $L = \lambda_{\max}(Q)$ и достигает каждой из них, когда $x - y$ — соответствующий собственный вектор. Отношение $\kappa = L/\mu$ показывает, насколько сильно расходятся две параболы, а в следующих разделах мы увидим, что оно определяет число итераций градиентного спуска.

\begin{selfcheckbox}
Примените неравенство Йенсена к $f(x) = x^2$ и случайной величине $X$ с конечной дисперсией. Что получится и чему равна разность его частей?
\selfcheckanswer Получится $(\mathbb E X)^2 \le \mathbb E X^2$. Разность равна $\mathbb E X^2 - (\mathbb E X)^2 = \operatorname{Var} X \ge 0$ и обращается в нуль, только если $X$ постоянна с вероятностью единица. Если $X = x$ с вероятностью $1-t$ и $X = y$ с вероятностью $t$, то $\operatorname{Var} X = t(1-t)(x - y)^2$: тождество для $x^2$ — это и есть формула для этой дисперсии.
\end{selfcheckbox}

% =====================================================================
\section{Градиентный спуск}\label{sec:gd}

Итерационный метод начинает с точки $x_0$ и повторяет
\handoutdisplay{
d_k = \mathcal D(x_k, f, \nabla f, \dots), \qquad \alpha_k = \mathcal S(x_k, d_k), \qquad x_{k+1} = x_k + \alpha_k d_k.
}
 Метод задаётся ответами на три вопроса: как выбрать направление $d_k$, как выбрать шаг $\alpha_k$ и когда остановиться. Градиентный спуск берёт $d_k = -\nabla f(x_k)$. В этом разделе шаг постоянный; бэктрекинг и критерии остановки — в следующем.

\begin{definitionbox}{Определение: Градиентный спуск.}
Пусть $f: \mathbb R^n \to \mathbb R$ дифференцируема и $\alpha > 0$. Градиентный спуск с постоянным шагом $\alpha$ из точки $x_0$ — это последовательность
\handoutdisplay{
x_{k+1} = x_k - \alpha\nabla f(x_k), \qquad k = 0, 1, 2, \dots
}
 Каждая итерация требует одного вычисления градиента.
\end{definitionbox}

\subsection{Почему антиградиент}\label{sub:gd-steepest}

\begin{theorembox}{Теорема: Наискорейшее возрастание в евклидовой норме.}
Пусть $f$ дифференцируема в точке $x$ и $\nabla f(x) \ne 0$. Тогда
\handoutdisplay{
\arg\max_{\|d\|_2 \le 1}\ \lim_{t \downarrow 0} \frac{f(x + td) - f(x)}{t} = \frac{\nabla f(x)}{\|\nabla f(x)\|_2},
}
 а наибольшая скорость возрастания равна $\|\nabla f(x)\|_2$.
\end{theorembox}

\begin{proofbox}
Из дифференцируемости следует, что $f(x + td) = f(x) + t\,\nabla f(x)^\top d + o(t)$, поэтому предел равен $\nabla f(x)^\top d$. По неравенству Коши — Буняковского — Шварца
\handoutdisplay{
\nabla f(x)^\top d \le \|\nabla f(x)\|_2\,\|d\|_2 \le \|\nabla f(x)\|_2,
}
 причём равенство достигается при $d = \nabla f(x)/\|\nabla f(x)\|_2$. Равенство в неравенстве Коши — Буняковского — Шварца возможно, только если $d$ получается из $\nabla f(x)$ умножением на неотрицательное число, а равенство во втором неравенстве требует $\|d\|_2 = 1$, так что эта точка максимума единственна.
\end{proofbox}

Заменив $d$ на $-d$, получаем: среди единичных векторов направление наискорейшего спуска — $-\nabla f(x)/\|\nabla f(x)\|_2$, а скорость изменения функции вдоль него равна $-\|\nabla f(x)\|_2$.

Слово \emph{наискорейший} относится к единичному шару евклидовой нормы. Для произвольной нормы наискорейшее направление единичной длины — решение задачи $\max_{\|d\| \le 1} \nabla f(x)^\top d = \|\nabla f(x)\|_*$, где $\|\cdot\|_*$ — двойственная норма, и ответ зависит от нормы. Для $\|d\|_M = \sqrt{d^\top M d}$ с $M \succ 0$ сделаем замену $u = M^{1/2} d$: задача превращается в $\max_{\|u\|_2 \le 1} \big(M^{-1/2}\nabla f(x)\big)^\top u$, теорема даёт $u \propto M^{-1/2}\nabla f(x)$, а значит, $d \propto M^{-1}\nabla f(x)$. В этой норме направление наискорейшего спуска идёт вдоль $-M^{-1}\nabla f(x)$; на этой идее основаны предобусловливание и натуральный градиент. Для $\|\cdot\|_\infty$ одно из направлений наискорейшего спуска — $-\operatorname{sign}\nabla f(x)$, где знак берётся покомпонентно.

\subsection{Лемма о спуске и безопасный шаг}\label{sub:gd-lemma}

\begin{theorembox}{Лемма: Лемма о спуске.}
Пусть $f: \mathbb R^n \to \mathbb R$ дифференцируема и её градиент липшицев с константой $L$. Тогда для любых $x, y$
\handoutdisplay{
f(y) \le f(x) + \nabla f(x)^\top (y - x) + \frac{L}{2}\|y - x\|^2.
}

\end{theorembox}

\begin{proofbox}
Положим $r(t) = f\big(x + t(y - x)\big)$ при $t \in [0, 1]$. Эта функция непрерывно дифференцируема, и $r'(t) = \nabla f\big(x + t(y - x)\big)^\top (y - x)$, поэтому
\handoutdisplay{
f(y) - f(x) = \int_0^1 r'(t)\,dt = \nabla f(x)^\top (y - x) + \int_0^1 \big(\nabla f(x + t(y - x)) - \nabla f(x)\big)^\top (y - x)\,dt.
}
 По неравенству Коши — Буняковского — Шварца и условию Липшица подынтегральное выражение не больше $\|\nabla f(x + t(y - x)) - \nabla f(x)\|\,\|y - x\| \le Lt\|y - x\|^2$, а $\int_0^1 Lt\,dt = \frac{L}{2}$.
\end{proofbox}

Та же оценка с модулями даёт $\big|f(y) - f(x) - \nabla f(x)^\top (y - x)\big| \le \frac{L}{2}\|y - x\|^2$: $L$-гладкая функция зажата между двумя параболами вокруг своей касательной плоскости. Градиентному спуску нужна только верхняя. Утверждение глобальное: одна константа $L$ оценивает ошибку линейной модели в любой точке и при любой длине шага.

\begin{theorembox}{Следствие: Убывание за один шаг.}
Пусть $f$ — $L$-гладкая функция, $\alpha > 0$ и $x^+ = x - \alpha\nabla f(x)$. Тогда
\handoutdisplay{
f(x^+) \le f(x) - \alpha\Big(1 - \frac{L\alpha}{2}\Big)\|\nabla f(x)\|^2.
}

\end{theorembox}

\begin{proofbox}
Подставим $y = x^+$ в лемму о спуске. Поскольку $y - x = -\alpha\nabla f(x)$, получаем
\handoutdisplay{
f(x^+) \le f(x) - \alpha\|\nabla f(x)\|^2 + \frac{L\alpha^2}{2}\|\nabla f(x)\|^2.
}

\end{proofbox}

Отсюда получаются две константы. При $0 < \alpha < 2/L$ множитель $1 - L\alpha/2$ положителен, поэтому каждый шаг из нестационарной точки строго уменьшает $f$. Гарантированное убывание $\alpha(1 - L\alpha/2)$ — вогнутая парабола по $\alpha$ с максимумом в точке $\alpha = 1/L$, где
\handoutdisplay{
f(x^+) \le f(x) - \frac{1}{2L}\|\nabla f(x)\|^2,
}
 и это лучшая гарантия, которую может дать такая оценка. При любом $\alpha \le 1/L$ выполнено $1 - L\alpha/2 \ge \tfrac{1}{2}$, а значит, $f(x^+) \le f(x) - \frac{\alpha}{2}\|\nabla f(x)\|^2$.

\begin{trapbox}
Значение $2/L$ — это граница, на которой кончается гарантия строгого убывания, а не доказательство того, что при более длинных шагах метод расходится: при $\alpha = 2/L$ оценка лишь исключает рост, при $\alpha > 2/L$ она его допускает, а конкретная функция может всё равно убывать. У квадратичной функции $2/L$ — точная граница, как показано в следующем подразделе.
\end{trapbox}

\subsection{Квадратичные функции и точное условие устойчивости}\label{sub:gd-quadratic}

Возьмём $f(x) = \tfrac{1}{2} x^\top Q x - b^\top x$ с симметричной матрицей $Q \succ 0$. Тогда $\nabla f(x) = Qx - b$, единственная точка минимума — $x^* = Q^{-1}b$, и $\nabla f(x) = Q(x - x^*)$. Для ошибки $e_k = x_k - x^*$ один шаг градиентного спуска даёт
\handoutdisplay{
e_{k+1} = e_k - \alpha Q e_k = (I - \alpha Q)\,e_k.
}
 Диагонализуем $Q = U\Lambda U^\top$, где $U$ ортогональна и $\Lambda = \operatorname{diag}(\lambda_1, \dots, \lambda_n)$, и перейдём к координатам в собственном базисе $z_k = U^\top e_k$. Координаты меняются независимо друг от друга:
\handoutdisplay{
z_{k+1,i} = (1 - \alpha\lambda_i)\,z_{k,i}, \qquad z_{k,i} = (1 - \alpha\lambda_i)^k z_{0,i},
}
 а $f(x_k) - f^* = \tfrac{1}{2} e_k^\top Q e_k = \tfrac{1}{2}\sum_i \lambda_i z_{k,i}^2$.

\begin{theorembox}{Теорема: Устойчивость на квадратичной функции.}
Пусть $f(x) = \tfrac{1}{2} x^\top Q x - b^\top x$ с симметричной $Q \succ 0$. Градиентный спуск с постоянным шагом $\alpha$ сходится к $x^*$ при любом начальном приближении в том и только в том случае, если
\handoutdisplay{
0 < \alpha < \frac{2}{\lambda_{\max}(Q)}.
}

\end{theorembox}

\begin{proofbox}
Из рекуррентного соотношения видно, что $z_{k,i} \to 0$ при любом $z_{0,i}$ в том и только в том случае, если $|1 - \alpha\lambda_i| < 1$, а $\|x_k - x^*\| = \|z_k\|$, так как $U$ ортогональна. Поэтому $x_k \to x^*$ из любой точки $x_0$ в том и только в том случае, если $|1 - \alpha\lambda_i| < 1$ для всех $i$. Поскольку $\lambda_i > 0$, это означает $0 < \alpha\lambda_i < 2$ при всех $i$, и определяющим оказывается наибольшее собственное значение.
\end{proofbox}

На границе $\alpha = 2/\lambda_{\max}$ множитель самой жёсткой координаты равен $-1$: эта координата без конца меняет знак, сохраняя модуль, а остальные по-прежнему уменьшаются. За границей она растёт в геометрической прогрессии. Между $1/\lambda_{\max}$ и $2/\lambda_{\max}$ её множитель отрицателен, но по модулю меньше $1$, и она стремится к нулю, чередуя знаки. Здесь $L = \lambda_{\max}(Q)$, так что безопасный диапазон $0 < \alpha < 2/L$ из леммы о спуске в точности совпадает с областью устойчивости: для квадратичных функций лемма неулучшаема.

\begin{examplebox}
В задаче A1 $Q = \operatorname{diag}(1;\ 10)$, поэтому $L = 10$ и граница равна $2/L = 0{,}2$. При $\alpha = 0{,}2$ множители равны $0{,}8$ и $-1$. Из $x_0 = (2;\ 2)$ жёсткая координата принимает попеременно значения $2$ и $-2$, и
\handoutdisplay{
f(x_k) = 20 + 2 \cdot 0{,}64^k:
}
 значения убывают к $20$ и никогда не опускаются ниже, потому что одна только жёсткая координата даёт вклад $\tfrac{1}{2} \cdot 10 \cdot 2^2 = 20$.
\end{examplebox}

\handoutfigure{0.9}{fig_stability}{Значения $f(x_k)$, логарифмическая шкала, для $f(x) = \tfrac{1}{2} x^\top Q x$ с $Q = \operatorname{diag}(1;\ 10)$, $x_0 = (2;\ 2)$ и четырьмя шагами из задачи A1. При шагах ниже границы $2/L = 0{,}2$ значения убывают в геометрической прогрессии, при $\alpha = 2/L$ стабилизируются на уровне $20$, а выше границы растут.}

\subsection{Обусловленность и оптимальный постоянный шаг}\label{sub:gd-conditioning}

Для квадратичной функции положим $\mu = \lambda_{\min}(Q)$ и $L = \lambda_{\max}(Q)$. За один шаг норма ошибки умножается не более чем на
\handoutdisplay{
\rho(\alpha) = \max_i |1 - \alpha\lambda_i| = \max\big(|1 - \alpha\mu|,\ |1 - \alpha L|\big),
}
 потому что $|1 - \alpha\lambda|$ выпукла по $\lambda$ и её максимум на $[\mu, L]$ достигается на конце отрезка. Вместе с $\|e_k\| = \|z_k\|$ это даёт
\handoutdisplay{
\|x_k - x^*\| \le \rho(\alpha)^k\,\|x_0 - x^*\|, \qquad f(x_k) - f^* \le \rho(\alpha)^{2k}\big(f(x_0) - f^*\big).
}


При $\alpha = 1/L$ множители $1 - \lambda_i/L$ лежат в $[0, 1 - 1/\kappa]$, поэтому $\rho(1/L) = 1 - 1/\kappa$. Так как $\ln\frac{1}{1 - 1/\kappa} \ge \frac{1}{\kappa}$, ошибка становится меньше $\varepsilon\|x_0 - x^*\|$, как только $k \ge \kappa\ln(1/\varepsilon)$: число итераций растёт пропорционально $\kappa$. Геометрически большое $\kappa$ — это длинная узкая долина; градиент направлен в основном поперёк неё, и шаг, достаточно короткий для крутого направления, почти не сдвигает точку вдоль пологого. Поэтому на квадратичной функции при $\alpha = 1/L$ невязка по функции за шаг умножается не более чем на $(1 - \mu/L)^2$. Это точнее общей гарантии $f(x_k) - f^* \le (1 - \mu/L)^k\big(f(x_0) - f^*\big)$, которая выполняется при $\alpha = 1/L$ для любой $\mu$-сильно выпуклой $L$-гладкой функции и доказывается в разделе «Скорость сходимости». В задачах A2 и A3 речь идёт именно об этой общей гарантии, поэтому множитель в них равен $1 - \mu/L$, а не $(1 - \mu/L)^2$.

\begin{theorembox}{Утверждение: Оптимальный постоянный шаг.}
Для квадратичной функции с $\mu = \lambda_{\min}(Q) > 0$ и $L = \lambda_{\max}(Q)$ минимум множителя $\rho(\alpha)$ по $\alpha > 0$ достигается при
\handoutdisplay{
\alpha^* = \frac{2}{L + \mu}, \qquad \rho(\alpha^*) = \frac{L - \mu}{L + \mu} = \frac{\kappa - 1}{\kappa + 1}.
}

\end{theorembox}

\begin{proofbox}
При $\alpha > 0$ имеем $|1 - \alpha\lambda| = \max(1 - \alpha\lambda,\ \alpha\lambda - 1)$; первое выражение максимально при $\lambda = \mu$, второе — при $\lambda = L$. Поэтому $\rho(\alpha) = \max(1 - \alpha\mu,\ \alpha L - 1)$ — максимум убывающей и возрастающей линейных функций от $\alpha$. Он минимален там, где они равны: из $1 - \alpha\mu = \alpha L - 1$ получаем $\alpha^* = 2/(L + \mu)$ и $\rho(\alpha^*) = 1 - \frac{2\mu}{L + \mu} = \frac{L - \mu}{L + \mu}$. При $\alpha < \alpha^*$ первая функция больше своего значения в $\alpha^*$, а при $\alpha > \alpha^*$ — вторая, так что любой другой шаг даёт больший множитель. Разделив числитель и знаменатель на $\mu$, получаем $(\kappa - 1)/(\kappa + 1)$.
\end{proofbox}

При $\alpha^*$ два крайних множителя равны $1 - \alpha^*\mu = \rho(\alpha^*)$ и $1 - \alpha^* L = -\rho(\alpha^*)$, то есть одинаковы по модулю. В задаче A3 значения $\mu = 1$ и $L = 50$ дают $\alpha^* = \tfrac{2}{51} \approx 0{,}039216$, а $\kappa = 50$ даёт $\rho(\alpha^*) = \tfrac{49}{51} \approx 0{,}960784$ против $1 - 1/\kappa = 0{,}98$ при $\alpha = 1/L$. Когда у $Q$ есть только собственные значения $\mu$ и $L$, как во всех двумерных примерах семинара, модуль каждой координаты за шаг умножается ровно на $\rho(\alpha^*)$, и $f(x_k) - f^* = \rho(\alpha^*)^{2k}\big(f(x_0) - f^*\big)$ выполняется как равенство. Для любой $\mu$-сильно выпуклой $L$-гладкой функции тот же шаг даёт ту же оценку для $\|x_k - x^*\|$ с тем же множителем, но доказательство длиннее.

\handoutfigure{0.95}{fig_conditioning}{Линии уровня $f(x) = \tfrac{1}{2} x^\top Q x$ и итерации градиентного спуска с $\alpha = 1/L$ из $x_0 = (1;\ 1)$. Слева: $Q = \operatorname{diag}(1;\ 2)$, $\kappa = 2$. Справа: $Q = \operatorname{diag}(1;\ 50)$, $\kappa = 50$. Пологая координата за шаг умножается на $1 - 1/\kappa$, то есть на $0{,}5$ слева и на $0{,}98$ справа, а чтобы добиться $f(x_k) \le 10^{-6}\,f(x_0)$, нужно 10 итераций слева и 245 справа.}

\subsection{Задачи семинара}\label{sub:gd-problems}

Задачи блока A с полными решениями.

\problemhead{A1}{Четыре шага на одной квадратичной функции}{У доски}
\textbf{Условие.} Минимизируйте $f(x) = \tfrac{1}{2} x^\top Q x$, где $Q = \operatorname{diag}(1;\ 10)$, градиентным спуском с постоянным шагом $\alpha$ из $x_0 = (2;\ 2)$. \textbf{а)} Найдите границу устойчивости: такое значение $\alpha_{\max}$, что метод сходится из любой начальной точки тогда и только тогда, когда $0 < \alpha < \alpha_{\max}$. \textbf{б)} Классифицируйте шаги $\alpha = 0{,}05$, $0{,}19$, $0{,}2$ и $0{,}21$ с помощью покоординатного рекуррентного соотношения $z_{k+1,i} = (1 - \alpha\lambda_i)\,z_{k,i}$: сходится ли метод и как именно?

\textbf{Решение.}
\begingroup\sloppy
\begin{enumerate}[label=\textbf{\arabic*.}, leftmargin=*]
  \item Здесь $x^* = 0$, $f^* = 0$, а $Q$ уже диагональна, так что собственные координаты совпадают с обычными: $z_k = x_k$, $\mu = \lambda_1 = 1$, $L = \lambda_2 = 10$, $\kappa = 10$. Шаг $x_{k+1} = x_k - \alpha Q x_k$ распадается на $z_{k+1,i} = (1 - \alpha\lambda_i)\,z_{k,i}$, и $f(x_k) = \tfrac{1}{2}\big(\lambda_1 z_{k,1}^2 + \lambda_2 z_{k,2}^2\big)$, поэтому $f(x_0) = 22$.
  \item Для сходимости из любой начальной точки нужно $|1 - \alpha\lambda_i| < 1$ при обоих $i$, то есть $0 < \alpha < 2/\lambda_i$. Пологая координата допускает $\alpha < 2/\mu = 2$, жёсткая — только $\alpha < 2/L = 0{,}2$. Границу задаёт жёсткое направление: $0 < \alpha < 0{,}2$.
  \item Множители для четырёх шагов:
\handoutdisplay{
\begin{array}{c|cc} \alpha & 1 - \alpha\lambda_1 & 1 - \alpha\lambda_2 \\ \hline 0{,}05 & 0{,}95 & 0{,}5 \\ 0{,}19 & 0{,}81 & -0{,}9 \\ 0{,}2 & 0{,}8 & -1 \\ 0{,}21 & 0{,}79 & -1{,}1 \end{array}
}

  \item $\alpha = 0{,}05$: оба множителя лежат в $(0, 1)$, поэтому обе координаты убывают, не меняя знака. Жёсткая уменьшается вдвое на каждом шаге, пологая теряет лишь двадцатую часть. После десяти шагов $f = 0{,}7170$, и почти всё это значение приходится на пологую координату. Сходится, но медленно.
  \item $\alpha = 0{,}19$: множитель жёсткой координаты $-0{,}9$ отрицателен, поэтому $z_{k,2}$ меняет знак на каждом шаге, но по модулю всё равно убывает. Сходится со сменой знака. Вначале преобладает жёсткая координата: на неё приходится $20$ из $f(x_0) = 22$, и здесь она по модулю умножается на $|-0{,}9|$, а при $\alpha = 0{,}05$ — на $0{,}5$, поэтому после десяти итераций этот шаг отстаёт: $f = 2{,}4611$ против $0{,}7170$. Но на длинной дистанции темп задаёт больший по модулю множитель, здесь $|-0{,}9|$ против $0{,}95$, и после пятидесяти итераций этот шаг далеко впереди: $f = 5{,}31 \cdot 10^{-4}$ против $1{,}18 \cdot 10^{-2}$.
  \item $\alpha = 0{,}2 = 2/L$: множитель жёсткой координаты равен ровно $-1$. Жёсткая координата без конца перескакивает между $+2$ и $-2$ и сама по себе на каждом шаге даёт в $f$ вклад $\tfrac{1}{2} \cdot 10 \cdot 2^2 = 20$. Пологая координата по-прежнему умножается на $0{,}8$, поэтому $f(x_k) = 20 + 2 \cdot 0{,}64^k$: значение убывает к $20$ и никогда не опускается ниже ($20{,}0231$ после десяти шагов). Ни сходимости к $x^*$, ни расходимости.
  \item $\alpha = 0{,}21$: $|-1{,}1| > 1$, поэтому $|z_{k,2}|$ растёт на каждом шаге, а знак чередуется, и $f(x_k) \to \infty$: $134{,}5679$ после десяти шагов, $2{,}76 \cdot 10^{5}$ после пятидесяти. Пологая координата сходится, но это не помогает. Расходится.
  \item Границу задаёт наибольшее собственное значение, а скорость — наименьшее: при любом $\alpha < 2/L$ множитель пологой координаты $1 - \alpha\mu > 1 - 2/\kappa$. При большом $\kappa$ приходится брать шаг, который почти не сдвигает пологую координату.
\end{enumerate}
\endgroup
\begin{answerbox}
Метод устойчив при $0 < \alpha < 2/L = 0{,}2$. $\alpha = 0{,}05$: сходится, медленно по пологой координате; $0{,}19$: сходится со сменой знака жёсткой координаты; $0{,}2$: жёсткая координата без конца перескакивает между $\pm2$, а $f(x_k)$ убывает к $20$, не опускаясь ниже; $0{,}21$: расходится.
\end{answerbox}

\problemhead{A2}{Гарантированное и фактическое число итераций}{Самостоятельно}
\textbf{Условие.} Запустите градиентный спуск для $f(x) = \tfrac{1}{2} x^\top Q x$, где $Q = \operatorname{diag}(1;\ 2)$, с шагом $\alpha = 1/L$ из $x_0 = (1;\ 1)$. Найдите \textbf{а)} $\kappa$; \textbf{б)} множитель $q$ в оценке из лекции $f(x_k) - f^* \le q^k\,\big(f(x_0) - f^*\big)$; \textbf{в)} наименьшее $k$, при котором эта оценка даёт $f(x_k) - f^* \le 10^{-6}\,\big(f(x_0) - f^*\big)$; \textbf{г)} наименьшее $k$, при котором итерации на самом деле удовлетворяют этому неравенству.

\textbf{Решение.}
\begingroup\sloppy
\begin{enumerate}[label=\textbf{\arabic*.}, leftmargin=*]
  \item $L = \lambda_{\max} = 2$, $\mu = \lambda_{\min} = 1$, поэтому $\kappa = 2$ и $\alpha = 1/L = 0{,}5$.
  \item Линейная скорость при $\alpha = 1/L$: $f(x_k) - f^* \le (1 - \mu/L)^k\,\big(f(x_0) - f^*\big)$, поэтому $q = 1 - \mu/L = 0{,}5$.
  \item Множитель в оценке опускается до $10^{-6}$, когда $0{,}5^k \le 10^{-6}$, то есть $k \ge \dfrac{\ln 10^{-6}}{\ln 0{,}5} = 19{,}93$. Округляем вверх: $20$ итераций.
  \item Теперь сами итерации; здесь $f^* = 0$. Множители $1 - \alpha\lambda_i$ равны $0{,}5$ для $\lambda = 1$ и $0$ для $\lambda = 2$: жёсткая координата обнуляется после первого же шага.
  \item Тогда $f(x_0) = \tfrac{1}{2}\big(1 + 2\big) = 1{,}5$ и $f(x_k) = \tfrac{1}{2} \cdot 0{,}5^{2k}$ при $k \ge 1$. Дальше невязка за итерацию умножается на $0{,}5^2$, а не на $0{,}5$.
  \item $f(x_k) \le 10^{-6}\,f(x_0)$ означает $0{,}5^{2k} \le 2 \cdot 1{,}5 \cdot 10^{-6}$, то есть $k \ge \dfrac{\ln\big(2 \cdot 1{,}5 \cdot 10^{-6}\big)}{2 \ln 0{,}5} = 9{,}17$. На деле нужно $10$ итераций.
  \item Оценка слабее, потому что в её доказательстве используются только неравенство из леммы о спуске и $\|\nabla f\|^2 \ge 2\mu\,(f - f^*)$, а они верны для любой функции с такими $\mu$ и $L$. На квадратичной функции при $\alpha = 1/L$ каждое координатное слагаемое невязки умножается на $(1 - \lambda_i/L)^2 \le (1 - \mu/L)^2$.
\end{enumerate}
\endgroup
\begin{answerbox}
$\kappa = 2$, $q = 0{,}5$; по оценке нужно $20$ итераций, на деле — $10$.
\end{answerbox}

\problemhead{A3}{Плохо обусловленная чаша и оптимальный шаг}{Самостоятельно}
\textbf{Условие.} Та же постановка, что в A2, но $Q = \operatorname{diag}(1;\ 50)$ и $x_0 = (1;\ 1)$. \textbf{а)–г)} Найдите те же четыре величины при $\alpha = 1/L$. \textbf{д)} Найдите значение оптимального постоянного шага $\alpha^* = 2/(L + \mu)$. \textbf{е)} При шаге $\alpha^*$ каждая координата ошибки $x_k - x^*$ на каждом шаге умножается на множитель, модуль которого не больше $\rho = (\kappa - 1)/(\kappa + 1)$. Сколько итераций нужно по получающейся отсюда оценке разности $f(x_k) - f^*$ для того же уменьшения — до доли $10^{-6}$ от начального значения?

\textbf{Решение.}
\begingroup\sloppy
\begin{enumerate}[label=\textbf{\arabic*.}, leftmargin=*]
  \item $L = 50$, $\mu = 1$, $\kappa = 50$, $\alpha = 1/L = 0{,}02$ и $q = 1 - \mu/L = 0{,}98$.
  \item Оценка: $k \ge \dfrac{\ln 10^{-6}}{\ln 0{,}98} = 683{,}84$, то есть $684$ итераций.
  \item Множители при $\alpha = 1/L$ равны $1 - \mu/L = 0{,}98$ и $1 - L/L = 0$. На жёсткую координату приходится $\tfrac{1}{2} \cdot 50$ из $f(x_0) = 25{,}5$, и после первого шага она обнуляется; дальше $f(x_k) = \tfrac{1}{2} \cdot 0{,}98^{2k}$ при $k \ge 1$.
  \item $f(x_k) \le 10^{-6}\,f(x_0)$ означает $0{,}98^{2k} \le 2 \cdot 25{,}5 \cdot 10^{-6}$, то есть $k \ge \dfrac{\ln\big(2 \cdot 25{,}5 \cdot 10^{-6}\big)}{2 \ln 0{,}98} = 244{,}6$. На деле нужно $245$ итераций, а по оценке — $684$.
  \item Оптимальный шаг: $\alpha^* = \dfrac{2}{L + \mu} = \tfrac{2}{51} \approx 0{,}039216$. Множители равны $1 - \alpha^*\mu = \tfrac{49}{51}$ и $1 - \alpha^* L = -\tfrac{49}{51}$: модули одинаковы, $\rho = \dfrac{\kappa - 1}{\kappa + 1} = \tfrac{49}{51} \approx 0{,}960784$.
  \item Невязка квадратична по координатам, поэтому $f(x_k) - f^* \le \rho^{2k}\,\big(f(x_0) - f^*\big)$. По оценке нужно $k \ge \dfrac{\ln 10^{-6}}{2 \ln \rho} = 172{,}67$, то есть $173$ итераций. Если взять $\rho^k$ вместо $\rho^{2k}$, получится $346$ — число итераций для расстояния $\|x_k - x^*\|$.
  \item Здесь обе координаты по модулю умножаются ровно на $\rho$, поэтому $f(x_k) = \rho^{2k} f(x_0)$ в точности, и на деле тоже нужно $173$ итераций. С шагом $\alpha^*$ оценка и фактическое число итераций совпадают; при $\alpha = 1/L$ они различаются почти втрое.
\end{enumerate}
\endgroup
\begin{answerbox}
$\kappa = 50$, $q = 0{,}98$; по оценке нужно $684$ итераций, на деле — $245$; $\alpha^* = \tfrac{2}{51} \approx 0{,}039216$, и оценка при $\alpha^*$ даёт $173$ — ровно столько, сколько нужно на деле.
\end{answerbox}

% =====================================================================
\section{Линейный поиск и критерии остановки}\label{sec:linesearch}

Для постоянного шага нужна константа $L$, которая редко известна, а шаг, безопасный там, где $f$ искривлена сильнее всего, в остальных местах слишком мал. Бэктрекинг выбирает шаг на каждой итерации по значениям $f$ вдоль луча.

\subsection{Бэктрекинг по Армихо}\label{sub:linesearch-armijo}

\begin{definitionbox}{Определение: Бэктрекинг по Армихо.}
Зафиксируем начальный шаг $\alpha_0 > 0$, коэффициент уменьшения $\beta \in (0, 1)$ и константу $c \in (0, 1)$, и пусть $d$ — направление спуска в точке $x$. Положим $\alpha = \alpha_0$ и будем повторять $\alpha \leftarrow \beta\alpha$ до тех пор, пока условие достаточного убывания (условие Армихо)
\handoutdisplay{
f(x + \alpha d) \le f(x) + c\,\alpha\,\nabla f(x)^\top d
}
 не выполнится. Принимается первый шаг вида $\alpha = \beta^j\alpha_0$, $j = 0, 1, 2, \dots$, который проходит проверку. Для градиентного спуска $d = -\nabla f(x)$, и правая часть равна $f(x) - c\,\alpha\|\nabla f(x)\|^2$. Обычно берут $\beta = \tfrac{1}{2}$ и $c = 10^{-4}$.
\end{definitionbox}

Вдоль луча положим $\varphi(\alpha) = f(x + \alpha d)$, так что $\varphi(0) = f(x)$ и $\varphi'(0) = \nabla f(x)^\top d < 0$. Прямая Армихо $\varphi(0) + c\alpha\varphi'(0)$ — это касательная, у которой наклон умножен на $c$, поэтому при $\alpha > 0$ она лежит выше касательной. Так как $\varphi(\alpha) = \varphi(0) + \alpha\varphi'(0) + o(\alpha)$, короткие шаги проходят проверку; длинный шаг не проходит, когда кривизна загибает $\varphi$ обратно выше этой прямой. Проверка отбрасывает шаги, убывание на которых составляет ничтожную долю того, что обещает наклон.

\begin{theorembox}{Теорема: Конечность бэктрекинга.}
Пусть $f$ — $L$-гладкая функция, $\nabla f(x) \ne 0$, $d = -\nabla f(x)$, $\alpha_0 > 0$ и $\beta, c \in (0, 1)$. Условие Армихо выполняется при всех $\alpha \in \big(0, 2(1-c)/L\big]$. Следовательно, бэктрекинг останавливается после конечного числа уменьшений, и принятый шаг удовлетворяет оценке
\handoutdisplay{
\alpha \ge \min\Big(\alpha_0,\ \frac{2\beta(1-c)}{L}\Big).
}

\end{theorembox}

\begin{proofbox}
По оценке убывания за один шаг $f\big(x - \alpha\nabla f(x)\big) \le f(x) - \alpha\big(1 - \frac{L\alpha}{2}\big)\|\nabla f(x)\|^2$. Если $\alpha \le 2(1-c)/L$, то $1 - L\alpha/2 \ge c$, и правая часть не больше $f(x) - c\alpha\|\nabla f(x)\|^2$: условие Армихо выполнено. Пробные шаги $\beta^j\alpha_0$ стремятся к $0$, поэтому при некотором конечном $j$ они попадают в $\big(0, 2(1-c)/L\big]$, и цикл остановится не позже этого. Если $\alpha_0$ проходит проверку, принимается шаг $\alpha_0$. Иначе принятому шагу $\alpha$ предшествует отвергнутый пробный шаг $\alpha/\beta$, каждый отвергнутый шаг больше $2(1-c)/L$, и поэтому $\alpha > 2\beta(1-c)/L$.
\end{proofbox}

Доказательство заодно оценивает объём работы: число уменьшений не превосходит наименьшего $j \ge 0$, при котором $\beta^j\alpha_0 \le 2(1-c)/L$, то есть $\big\lceil \log_{1/\beta}\big(\alpha_0 L/(2(1-c))\big) \big\rceil$ при $\alpha_0 > 2(1-c)/L$. Цикл может остановиться раньше, потому что вдоль конкретного направления условие может выполняться и за пределами $2(1-c)/L$. Для квадратичной функции диапазон известен точно: обозначим $g = \nabla f(x)$, тогда
\handoutdisplay{
f(x - \alpha g) = f(x) - \alpha\,g^\top g + \frac{\alpha^2}{2}\,g^\top Q g,
}
 поэтому условие Армихо выполняется тогда и только тогда, когда $\alpha \le 2(1-c)\,\dfrac{g^\top g}{g^\top Q g}$, и эта величина не меньше $2(1-c)/L$, потому что $g^\top Q g \le L\,g^\top g$.

\begin{examplebox}
В задаче B1 $Q = \operatorname{diag}(1;\ 10)$, $x_0 = (2;\ 2)$, $L = 10$ и $c = 10^{-4}$. Теорема гарантирует принятие шага при $\alpha \le 2(1-c)/L = 0{,}19998$, а вдоль $d = -\nabla f(x_0)$ точный диапазон — $\alpha \le 0{,}20178$. Если делить шаг пополам, начиная с $\alpha_0 = 1$ первый пробный шаг, попадающий в оба диапазона, равен $0{,}125$; он получается после 3 делений и лежит между $1/L = 0{,}1$ и $2/L = 0{,}2$.
\end{examplebox}

Без $L$-гладкости бэктрекинг всё равно завершается для любого направления спуска $d$ дифференцируемой функции $f$, потому что $f(x + \alpha d) - f(x) - c\alpha\nabla f(x)^\top d = (1 - c)\,\alpha\,\nabla f(x)^\top d + o(\alpha)$ отрицательно при малых $\alpha$. Пропадает только оценка снизу на принятый шаг.

\begin{trapbox}
Если $d$ не является направлением спуска, короткие пробные шаги не проходят проверку, и цикл может никогда не закончиться. Поэтому в реализации проверяют $\nabla f(x)^\top d < 0$ и ограничивают число уменьшений.
\end{trapbox}

\subsection{Критерии остановки}\label{sub:linesearch-stopping}

В лекции используются три масштабированных критерия с допусками $\varepsilon_g$, $\varepsilon_x$ и $\varepsilon_f$:
\handoutdisplay{
\begin{gathered}
\|\nabla f(x_k)\| \le \varepsilon_g \max\big(1, |f(x_k)|\big), \\[6pt]
\|x_{k+1} - x_k\| \le \varepsilon_x \max\big(1, \|x_k\|\big), \\[6pt]
|f(x_{k+1}) - f(x_k)| \le \varepsilon_f \max\big(1, |f(x_k)|\big).
\end{gathered}
}
 Множитель $\max(1, \cdot)$ делает проверку относительной, когда масштаб велик, и абсолютной, когда он меньше $1$, так что от величины, близкой к нулю, не требуется недостижимая относительная точность.

В коде допуск часто разбивают на абсолютную и относительную части, $\text{atol} + \text{rtol}\cdot\max(1, s)$, где $s$ — масштаб: $|f(x_k)|$ или $\|x_k\|$. В задаче B3 используется именно такая форма с $\text{atol} = 10^{-8}$ и $\text{rtol} = 10^{-6}$.

Критерии измеряют разное. Критерии по изменению следят за движением метода: для градиентного спуска $\|x_{k+1} - x_k\| = \alpha_k\|\nabla f(x_k)\|$ и $f(x_k) - f(x_{k+1}) \approx \alpha_k\|\nabla f(x_k)\|^2$, и обе величины малы, как только мал шаг, как бы далеко точка ни была от решения. Малые изменения могут означать сходимость, крошечный шаг, плохое масштабирование или застой из-за арифметики с плавающей точкой. Критерий по градиенту смотрит на саму точку. Малый градиент — сертификат стационарности первого порядка; при условии Поляка–Лоясевича (PL) он оценивает ошибку по $f$, а для $\mu$-сильно выпуклой $f$ — ещё и расстояние до $x^*$:
\handoutdisplay{
f(x) - f^* \le \frac{1}{2\mu}\|\nabla f(x)\|^2, \qquad \|x - x^*\| \le \frac{1}{\mu}\|\nabla f(x)\|.
}
 Первое неравенство — это переписанное условие PL. Для второго запишем неравенство сильной выпуклости первого порядка в точках $x$ и $x^*$ и сложим: получится $\big(\nabla f(x) - \nabla f(x^*)\big)^\top (x - x^*) \ge \mu\|x - x^*\|^2$; остаётся воспользоваться тем, что $\nabla f(x^*) = 0$, и неравенством Коши — Буняковского — Шварца.

\begin{notebox}{Комбинированное правило}
Сообщайте об успехе и останавливайтесь, когда выполнен критерий по градиенту с допуском, умноженным на $\max(1, |f(x_k)|)$, как в лекции. Критерии по изменению считайте лишь сигналом застоя: если они выполнены, а критерий по градиенту нет, остановитесь и сообщите, что метод застрял, а не сошёлся. Всегда добавляйте ограничение на число итераций и вычислений функции и останавливайтесь, если $f$ или $\nabla f$ перестают быть конечными.
\end{notebox}

\subsection{Критерий по изменению, срабатывающий слишком рано}\label{sub:linesearch-failure}

В задаче B3 такой сбой построен намеренно. Возьмём $f(x) = \tfrac{1}{2} x^\top Q x$ с $Q = \operatorname{diag}(1;\ 50)$, так что $x^* = 0$ и $f^* = 0$, начнём с $x_0 = (1;\ 1)$, где $f(x_0) = 25{,}5$, и сделаем один шаг градиентного спуска с заданным шагом $\alpha = 10^{-8}$. Получаем $x_1 = (0{,}99999999,\ 0{,}9999995)$ и $f(x_1) = 25{,}49997499$. Три критерия с $\text{atol} = 10^{-8}$ и $\text{rtol} = 10^{-6}$ дают следующее.

\begin{center}\small
\begin{tabular}{@{}>{\raggedright\arraybackslash\hspace{0pt}}p{0.17\linewidth}rll@{}}
\toprule
Критерий & Значение & Порог & Срабатывает \\
\midrule
изменение функции $|f(x_1) - f(x_0)|$ & $2{,}501 \cdot 10^{-5}$ & $\text{atol} + \text{rtol}\max(1, |f(x_0)|) = 2{,}551 \cdot 10^{-5}$ & да \\
шаг $\|x_1 - x_0\|$ & $5{,}001 \cdot 10^{-7}$ & $\text{atol} + \text{rtol}\max(1, \|x_0\|) = 1{,}424 \cdot 10^{-6}$ & да \\
градиент $\|\nabla f(x_1)\|$ & $50{,}01$ & $\text{atol} + \text{rtol}\max(1, |f(x_1)|) = 2{,}551 \cdot 10^{-5}$ & нет \\
\bottomrule
\end{tabular}
\end{center}

Оба критерия по изменению срабатывают после одного шага, хотя $f$ уменьшилась меньше чем на миллионную долю своего значения, а точка почти не приблизилась к $x^*$. Только критерий по градиенту не даёт остановиться. Критерии по изменению срабатывают потому, что их левые части пропорциональны шагу: $\|x_1 - x_0\| = \alpha\|\nabla f(x_0)\|$ и $|f(x_1) - f(x_0)| \approx \alpha\|\nabla f(x_0)\|^2$, где $\|\nabla f(x_0)\| = 50{,}01$, и крошечный шаг делает обе величины крошечными. В критерии по градиенту множителя $\alpha$ нет — он измеряет градиент в самой точке $x_1$, а после крошечного шага градиент почти такой же, как в $x_0$: $\|\nabla f(x_1)\| \approx \|\nabla f(x_0)\|$.

\subsection{Задачи семинара}\label{sub:linesearch-problems}

Задачи блока B с полными решениями.

\problemhead{B1}{Бэктрекинг по Армихо вручную}{У доски}
\textbf{Условие.} Пусть $f(x) = \tfrac{1}{2} x^\top Qx$, где $Q = \operatorname{diag}(1, 10)$, и $x_0 = (2;\ 2)$. Выполните бэктрекинг по Армихо вдоль $d = -\nabla f(x_0)$ с $\alpha_0 = 1$, $\beta = 0{,}5$ и $c = 10^{-4}$: сначала попробуйте $\alpha = \alpha_0$ и заменяйте $\alpha \leftarrow \beta\alpha$, пока $f(x_0 + \alpha d) > f(x_0) + c\alpha\nabla f(x_0)^\top d$. Сколько раз шаг делится пополам, прежде чем условие выполнится, какой шаг принимается и как он соотносится с $1/L$ и $2/L$?

\textbf{Решение.}
\begingroup\sloppy
\begin{enumerate}[label=\textbf{\arabic*.}, leftmargin=*]
  \item Данные в точке $x_0$: $f(x_0) = 22$, $\nabla f(x_0) = Qx_0 = (2;\ 20)$ и $\nabla f(x_0)^\top d = -\|\nabla f(x_0)\|^2 = -404$. Здесь $L = 10$, $1/L = 0{,}1$ и $2/L = 0{,}2$.
  \item Вдоль луча у каждой координаты свой множитель: $x_0 + \alpha d = \big(2(1 - \alpha),\ 2(1 - 10\alpha)\big)$. Условие принимает вид $f(x_0 + \alpha d) \le 22 - 10^{-4}\cdot404\,\alpha$.
  \item Попытки с делением шага пополам, начиная с $\alpha_0$:
\handoutdisplay{
\begin{array}{r|r|r|c} \alpha & f(x_0 + \alpha d) & f(x_0) + c\alpha\nabla f(x_0)^\top d & \text{выполнено} \\ \hline 1 & 1620 & 21{,}9596 & \text{нет} \\ 0{,}5 & 320{,}5 & 21{,}9798 & \text{нет} \\ 0{,}25 & 46{,}125 & 21{,}9899 & \text{нет} \\ 0{,}125 & 2{,}78125 & 21{,}99495 & \text{да} \end{array}
}

  \item Первая попытка, прошедшая проверку, — $\alpha = 0{,}125$ после $3$ делений пополам, и на ней цикл останавливается. Этот шаг лежит между $1/L = 0{,}1$ и $2/L = 0{,}2$, то есть внутри области устойчивости из блока A.
  \item Почему цикл обязательно остановится: лемма о спуске даёт $f(x_0 + \alpha d) \le f(x_0) - \alpha(1 - L\alpha/2)\|\nabla f(x_0)\|^2$, поэтому условие выполнено при всех $\alpha \le 2(1-c)/L = 0{,}19998$. Каждая отвергнутая попытка больше этой границы, так что принятый шаг не меньше этой границы, умноженной на $\beta$.
  \item Для квадратичной функции неравенство Армихо решается точно. Обозначим $g = \nabla f(x_0)$; тогда $f(x_0 - \alpha g) = f(x_0) - \alpha\|g\|^2 + \tfrac{1}{2}\alpha^2 g^\top Qg$, и условие выполнено при $\alpha \le 2(1-c)\|g\|^2/(g^\top Qg) = 0{,}20178$. Это чуть больше $0{,}19998$, потому что кривизна вдоль $g$ меньше $L$. Минимум вдоль луча достигается при шаге $\|g\|^2/(g^\top Qg) = 0{,}1009$, но бэктрекинг по Армихо его не ищет.
\end{enumerate}
\endgroup
\begin{answerbox}
Число делений пополам: $3$; принятый шаг $\alpha = 0{,}125$ лежит между $1/L = 0{,}1$ и $2/L = 0{,}2$; условие выполнено при всех $\alpha \le 2(1-c)/L = 0{,}19998$.
\end{answerbox}

\problemhead{B2}{Бэктрекинг на более жёсткой квадратичной функции}{Самостоятельно}
\textbf{Условие.} Пусть $f(x) = \tfrac{1}{2} x^\top Qx$, где $Q = \operatorname{diag}(1, 50)$, и $x_0 = (2;\ 2)$. Выполните тот же бэктрекинг, что в B1, вдоль $d = -\nabla f(x_0)$ с $\alpha_0 = 1$, $\beta = 0{,}5$ и $c = 10^{-4}$: сначала попробуйте $\alpha = \alpha_0$ и заменяйте $\alpha \leftarrow \beta\alpha$, пока $f(x_0 + \alpha d) > f(x_0) + c\alpha\nabla f(x_0)^\top d$. Сколько раз шаг делится пополам, какой шаг принимается и где он лежит относительно $1/L$ и $2/L$?

\textbf{Решение.}
\begingroup\sloppy
\begin{enumerate}[label=\textbf{\arabic*.}, leftmargin=*]
  \item Данные: $f(x_0) = 102$, $\nabla f(x_0) = (2;\ 100)$, $\|\nabla f(x_0)\|^2 = 10004$; $L = 50$, $1/L = 0{,}02$, $2/L = 0{,}04$. Так как $x_0 + \alpha d = \big(2(1 - \alpha),\ 2(1 - 50\alpha)\big)$, условие имеет вид $f(x_0 + \alpha d) \le 102 - 10^{-4}\cdot10004\,\alpha$.
  \item Таблица попыток: шаг делится пополам, начиная с $\alpha_0$:
\handoutdisplay{
\begin{array}{r|r|r|c} \alpha & f(x_0 + \alpha d) & f(x_0) + c\alpha\nabla f(x_0)^\top d & \text{выполнено} \\ \hline 1 & 240100 & 100{,}9996 & \text{нет} \\ 0{,}5 & 57600{,}5 & 101{,}4998 & \text{нет} \\ 0{,}25 & 13226{,}125 & 101{,}7499 & \text{нет} \\ 0{,}125 & 2757{,}78125 & 101{,}87495 & \text{нет} \\ 0{,}0625 & 453{,}3203125 & 101{,}937475 & \text{нет} \\ 0{,}03125 & 33{,}517578125 & 101{,}9687375 & \text{да} \end{array}
}

  \item Условие Армихо впервые выполняется при $\alpha = 0{,}03125$, после $5$ делений пополам. Как и в B1, этот шаг лежит между $1/L = 0{,}02$ и $2/L = 0{,}04$.
  \item Гарантия: любой шаг $\alpha \le 2(1-c)/L = 0{,}039996$ проходит проверку, а эта граница чуть меньше $2/L$. Попытки пересекают эту границу между $0{,}0625$ и $0{,}03125$, и принятый шаг не меньше этой границы, умноженной на $\beta$.
  \item Точный порог для этой квадратичной функции: $2(1-c)\|g\|^2/(g^\top Qg) = 0{,}040012$ при $g = \nabla f(x_0)$, а это чуть больше $2/L$. Условие Армихо проверяет убывание за один шаг, а не устойчивость: попытка чуть больше $2/L$ здесь прошла бы проверку, хотя градиентный спуск с таким постоянным шагом расходится.
  \item По сравнению с B1 константа $L$ выросла с $10$ до $50$, и бэктрекингу понадобилось $5$ делений пополам вместо $3$. В обеих задачах принятый шаг попадает между $1/L$ и $2/L$, хотя $L$ бэктрекинг ни разу не использовал.
\end{enumerate}
\endgroup
\begin{answerbox}
Число делений пополам: $5$; принятый шаг $\alpha = 0{,}03125$ лежит между $1/L = 0{,}02$ и $2/L = 0{,}04$ и меньше гарантированной границы $2(1-c)/L = 0{,}039996$.
\end{answerbox}

\problemhead{B3}{Три критерия остановки после крошечного шага}{Самостоятельно}
\textbf{Условие.} Пусть $f(x) = \tfrac{1}{2} x^\top Qx$, где $Q = \operatorname{diag}(1, 50)$, и $x_0 = (1;\ 1)$. Сделайте один градиентный шаг $x_1 = x_0 - \alpha\nabla f(x_0)$ с заданным шагом $\alpha = 10^{-8}$. При $\mathrm{atol} = 10^{-8}$, $\mathrm{rtol} = 10^{-6}$ и евклидовых нормах проверьте три критерия с учётом масштаба
\handoutdisplay{
\begin{aligned} \text{градиент:}\quad & \|\nabla f(x_1)\| \le \mathrm{atol} + \mathrm{rtol}\cdot\max(1, |f(x_1)|), \\ \text{шаг:}\quad & \|x_1 - x_0\| \le \mathrm{atol} + \mathrm{rtol}\cdot\max(1, \|x_0\|), \\ \text{функция:}\quad & |f(x_1) - f(x_0)| \le \mathrm{atol} + \mathrm{rtol}\cdot\max(1, |f(x_0)|). \end{aligned}
}
 Критерий срабатывает, когда его неравенство выполнено. Найдите $|f(x_1) - f(x_0)|$, $\|\nabla f(x_1)\|$ и число сработавших критериев. Близка ли точка $x_1$ к точке минимума? Затем сформулируйте комбинированный критерий остановки и укажите, на какую величину он нормирован.

\textbf{Решение.}
\begingroup\sloppy
\begin{enumerate}[label=\textbf{\arabic*.}, leftmargin=*]
  \item Один шаг: $\nabla f(x_0) = Qx_0 = (1, 50)$, $\|\nabla f(x_0)\| = 50{,}01$ и $x_1 = x_0 - \alpha\nabla f(x_0) = (0{,}99999999,\ 0{,}9999995)$.
  \item Критерий по функции: $f(x_0) = 25{,}5$ и $f(x_1) = 25{,}49997499$, поэтому $|f(x_1) - f(x_0)| = 2{,}501 \cdot 10^{-5}$, что совпадает с первым приближением $\alpha\|\nabla f(x_0)\|^2$. Порог равен $\mathrm{atol} + \mathrm{rtol}\cdot25{,}5 = 2{,}551 \cdot 10^{-5}$: критерий срабатывает.
  \item Критерий по шагу: $\|x_1 - x_0\| = \alpha\|\nabla f(x_0)\| = 5{,}001 \cdot 10^{-7}$ при пороге $\mathrm{atol} + \mathrm{rtol}\cdot\max(1, \|x_0\|) = \mathrm{atol} + \mathrm{rtol}\sqrt2 = 1{,}424 \cdot 10^{-6}$: он тоже срабатывает.
  \item Критерий по градиенту: $\|\nabla f(x_1)\| = 50{,}01$, почти как в точке $x_0$, а порог равен $\mathrm{atol} + \mathrm{rtol}\cdot\max(1, |f(x_1)|) = 2{,}551 \cdot 10^{-5}$. Этот критерий не срабатывает, так что из $3$ критериев срабатывают $2$.
  \item Точка минимума — $x^* = 0$, и в ней $f^* = 0$, а у нас $\|x_1 - x^*\| \approx \sqrt2$ и $f(x_1) \approx f(x_0)$. Критерии по изменениям сработали потому, что оба изменения уменьшаются вместе с $\alpha$ (в первом приближении они ему пропорциональны): они измеряют шаг, а не расстояние до $x^*$.
  \item Общее правило: объявляйте сходимость, только если $\|\nabla f(x_k)\| \le \mathrm{atol} + \mathrm{rtol}\cdot\max(1, |f(x_k)|)$, то есть по норме градиента с масштабом $\max(1, |f(x_k)|)$. Если срабатывают только критерии по изменениям, останавливайтесь с предупреждением о застое и не выдавайте точку за решение; всегда ограничивайте число итераций и останавливайтесь при бесконечных или нечисловых значениях.
\end{enumerate}
\endgroup
\begin{answerbox}
$|f(x_1) - f(x_0)| = 2{,}501 \cdot 10^{-5}$ и $\|x_1 - x_0\| = 5{,}001 \cdot 10^{-7}$ меньше своих порогов, а $\|\nabla f(x_1)\| = 50{,}01$ — нет: $2$ критерия срабатывают далеко от $x^* = 0$. Останавливайтесь по норме градиента с масштабом $\max(1, |f(x_k)|)$, а критерии по изменениям считайте только сигналом застоя и ограничивайте число итераций.
\end{answerbox}

% =====================================================================
\section{Скорость сходимости}\label{sec:rates}

В лекции доказаны три оценки для градиентного спуска с шагом $\alpha = 1/L$: для выпуклых функций невязка по функции имеет порядок $1/k$, а для сильно выпуклых функций и, в более общем случае, при условии PL она на каждом шаге умножается на $1 - \mu/L$. Все три доказательства начинаются с неравенства для шага спуска из раздела о градиентном спуске: для $L$-гладкой $f$ и $x_{k+1} = x_k - \frac1L \nabla f(x_k)$ лемма о спуске даёт
\handoutdisplay{
f(x_{k+1}) \le f(x_k) - \frac{1}{2L}\|\nabla f(x_k)\|^2.
}
 В частности, $f(x_k)$ никогда не возрастает. Всюду ниже $f^* = \inf f$, а $x^*$ — точка минимума, если она существует.

\subsection{Выпуклый случай}\label{sub:rates-convex}

\begin{theorembox}{Теорема: Скорость сходимости для выпуклых функций.}
Пусть $f$ — выпуклая $L$-гладкая функция, $x^*$ — точка минимума, $f^* = f(x^*)$. Тогда для градиентного спуска с $\alpha = 1/L$ выполнено
\handoutdisplay{
f(x_k) - f^* \le \frac{L\|x_0 - x^*\|^2}{2k} \qquad \text{для всех } k \ge 1.
}

\end{theorembox}

\begin{proofbox}
Обозначим $g_k = \nabla f(x_k)$. По выпуклости касательная плоскость в точке $x_k$ лежит под графиком, поэтому $f^* \ge f(x_k) + g_k^\top(x^* - x_k)$, то есть $f(x_k) \le f^* + g_k^\top(x_k - x^*)$. Подставив это в неравенство для шага спуска, получаем
\handoutdisplay{
f(x_{k+1}) - f^* \le g_k^\top(x_k - x^*) - \frac{1}{2L}\|g_k\|^2.
}
 Правая часть — разность квадратов расстояний. Так как $x_{k+1} - x^* = (x_k - x^*) - \frac1L g_k$, раскрыв квадрат, получаем тождество трёх точек
\handoutdisplay{
\|x_{k+1} - x^*\|^2 = \|x_k - x^*\|^2 - \frac2L\, g_k^\top(x_k - x^*) + \frac{1}{L^2}\|g_k\|^2,
}
 а умножив это тождество на $L/2$, превращаем предыдущее неравенство в ключевое:
\handoutdisplay{
f(x_{k+1}) - f^* \le \frac L2\Big(\|x_k - x^*\|^2 - \|x_{k+1} - x^*\|^2\Big).
}
 Просуммируем его по шагам $0, 1, \dots, k-1$: правая часть сокращается телескопически:
\handoutdisplay{
\sum_{j=1}^{k}\big(f(x_j) - f^*\big) \le \frac L2\Big(\|x_0 - x^*\|^2 - \|x_k - x^*\|^2\Big) \le \frac L2\|x_0 - x^*\|^2.
}
 Наконец, последовательность $f(x_j)$ не возрастает, поэтому каждое из $k$ слагаемых слева не меньше $f(x_k) - f^*$. Отсюда $k\big(f(x_k) - f^*\big) \le \frac L2\|x_0 - x^*\|^2$.
\end{proofbox}

У ключевого неравенства есть и второе следствие. Его левая часть неотрицательна, поэтому $\|x_{k+1} - x^*\| \le \|x_k - x^*\|$: итерации никогда не удаляются ни от одной точки минимума.

\begin{whybox}
Каждый шаг оплачивает свою невязку уменьшением квадрата расстояния до $x^*$, а весь бюджет равен $\|x_0 - x^*\|^2$. Поэтому в среднем по первым $k$ шагам невязки малы, а раз они только убывают, последняя не больше среднего.
\end{whybox}

Скорость \textbf{сублинейная}: чтобы уменьшить оценку вдвое, нужно вдвое больше шагов, а невязка меньше $\varepsilon$ гарантирована только после $L\|x_0 - x^*\|^2/(2\varepsilon)$ шагов.

\begin{selfcheckbox}
Теорема верна для любой точки минимума $x^*$. Если точки минимума образуют целое множество, какая из них даёт лучшую оценку?
\selfcheckanswer Ближайшая к $x_0$. Во всех точках минимума функция принимает одно и то же значение $f^*$, а больше ничего об $x^*$ доказательство не использует, поэтому $\|x_0 - x^*\|$ можно заменить расстоянием от $x_0$ до множества точек минимума.
\end{selfcheckbox}

\subsection{Сильная выпуклость}\label{sub:rates-strong}

Дифференцируемая функция $f$ называется $\mu$-\textbf{сильно выпуклой}, $\mu > 0$, если
\handoutdisplay{
f(y) \ge f(x) + \nabla f(x)^\top (y - x) + \frac\mu2\|y - x\|^2 \qquad \text{для всех } x, y.
}
 Для $L$-гладкой $f$ лемма о спуске добавляет оценку сверху, и выражение $f(y) - f(x) - \nabla f(x)^\top(y - x)$ оказывается зажатым между $\frac\mu2\|y - x\|^2$ и $\frac L2\|y - x\|^2$, поэтому $\mu \le L$ и $\kappa = L/\mu \ge 1$.

\begin{theorembox}{Лемма: Из сильной выпуклости следует условие PL.}
Если $f$ — $\mu$-сильно выпуклая функция, то для любого $x$
\handoutdisplay{
\frac{1}{2}\|\nabla f(x)\|^2 \ge \mu\big(f(x) - f^*\big).
}

\end{theorembox}

\begin{proofbox}
Зафиксируем $x$ и обозначим $g = \nabla f(x)$. Правая часть определения — квадратичная функция от $y$:
\handoutdisplay{
q(y) = f(x) + g^\top (y - x) + \frac\mu2\|y - x\|^2.
}
 Её градиент $g + \mu(y - x)$ обращается в нуль при $y = x - g/\mu$, поэтому её наименьшее значение равно
\handoutdisplay{
q\Big(x - \frac g\mu\Big) = f(x) - \frac1\mu\|g\|^2 + \frac{1}{2\mu}\|g\|^2 = f(x) - \frac{1}{2\mu}\|g\|^2.
}
 Значит, $f(y) \ge q(y) \ge f(x) - \frac{1}{2\mu}\|g\|^2$ для любого $y$. Переходя к инфимуму по $y$, получаем $f^* \ge f(x) - \frac{1}{2\mu}\|g\|^2$; остаётся умножить на $\mu$ и перенести слагаемые.
\end{proofbox}

\begin{theorembox}{Теорема: Линейная скорость при сильной выпуклости.}
Пусть $f$ — $\mu$-сильно выпуклая и $L$-гладкая функция. Тогда для градиентного спуска с $\alpha = 1/L$ выполнено
\handoutdisplay{
f(x_{k+1}) - f^* \le \Big(1 - \frac\mu L\Big)\big(f(x_k) - f^*\big), \qquad f(x_k) - f^* \le \Big(1 - \frac\mu L\Big)^k\big(f(x_0) - f^*\big).
}

\end{theorembox}

\begin{proofbox}
Вычтем $f^*$ из обеих частей неравенства для шага спуска и оценим слагаемое с градиентом снизу по лемме: $\frac{1}{2L}\|\nabla f(x_k)\|^2 \ge \frac\mu L\big(f(x_k) - f^*\big)$. Получаем
\handoutdisplay{
f(x_{k+1}) - f^* \le f(x_k) - f^* - \frac{1}{2L}\|\nabla f(x_k)\|^2 \le \Big(1 - \frac\mu L\Big)\big(f(x_k) - f^*\big).
}
 Применив это неравенство $k$ раз, получаем второе.
\end{proofbox}

Множитель $1 - \mu/L = 1 - 1/\kappa$ лежит в $[0, 1)$. Оценка вида $C q^k$ с $q < 1$ называется \textbf{линейной скоростью сходимости}: в логарифмическом масштабе оценка убывает по прямой.

\begin{whybox}
Сильная выпуклость не даёт градиенту быть малым, пока невязка велика. Поэтому с шагом $1/L$ каждая итерация убирает по меньшей мере долю $\mu/L$ оставшейся невязки.
\end{whybox}

\begin{selfcheckbox}
Возьмём $f(x) = \frac{1}{2}x^\top Q x$ с $\kappa = 1$. Что даёт теорема и что на самом деле делает градиентный спуск с $\alpha = 1/L$?
\selfcheckanswer Множитель равен $1 - \mu/L = 0$, поэтому невязка обращается в нуль за один шаг. Действительно, $\kappa = 1$ означает, что $Q = LI$, и $x_1 = x_0 - \frac1L \cdot L x_0 = 0 = x^*$.
\end{selfcheckbox}

\subsection{Условие Поляка–Лоясевича}\label{sub:rates-pl}

Доказательство выше использовало сильную выпуклость только один раз — через неравенство из леммы. Это неравенство и есть условие PL, и его одного достаточно.

\begin{theorembox}{Теорема: Линейная скорость при условии PL.}
Пусть $f$ — $L$-гладкая функция с $f^* = \inf f > -\infty$, удовлетворяющая условию Поляка–Лоясевича с некоторым $\mu > 0$:
\handoutdisplay{
\frac{1}{2}\|\nabla f(x)\|^2 \ge \mu\big(f(x) - f^*\big) \qquad \text{для всех } x.
}
 Тогда для градиентного спуска с $\alpha = 1/L$ выполнено $f(x_k) - f^* \le (1 - \mu/L)^k\big(f(x_0) - f^*\big)$.
\end{theorembox}

\begin{proofbox}
Неравенство для шага спуска даёт $f(x_{k+1}) \le f(x_k) - \frac{1}{2L}\|\nabla f(x_k)\|^2$, а условие PL оценивает вычитаемое снизу величиной $\frac\mu L\big(f(x_k) - f^*\big)$. Вычитая $f^*$, получаем $f(x_{k+1}) - f^* \le (1 - \mu/L)\big(f(x_k) - f^*\big)$, и остаётся применить это неравенство шаг за шагом. Выпуклость нигде не использовалась.
\end{proofbox}

Условие PL означает, что градиент велик везде, где велика невязка. В стационарной точке левая часть равна $0$, поэтому $f(x) = f^*$: \textbf{любая стационарная точка — точка глобального минимума}. Обычно для этого и нужна выпуклость, а PL даёт то же самое без неё.

\begin{examplebox}
Функция $f(x) = x^2 + 3\sin^2 x$ на $\mathbb R$ удовлетворяет условию PL, но не выпукла. Она не выпукла, потому что $f''(x) = 2 + 6\cos 2x$ при $x = \pi/2$ равна $2 - 6 < 0$. Она $L$-гладкая с $L = 8$, так как $|f''| \le 8$. Её единственная стационарная точка — точка минимума $0$: при $0 < x < \tfrac{3}{2}$ оба слагаемых в $f'(x) = 2x + 3\sin 2x$ положительны, при $x \ge \tfrac{3}{2}$ имеем $2x \ge 3 \ge -3\sin 2x$, причём оба неравенства не обращаются в равенства одновременно, а $f'$ нечётна. Поэтому отношение $\frac{1}{2}f'(x)^2/f(x)$ непрерывно и положительно при $x \ne 0$. При $x \to 0$ оно стремится к $8$, поскольку там $f(x) \approx 4x^2$ и $f'(x) \approx 8x$, а при $|x| \to \infty$ — к $2$. Значит, его инфимум положителен, и условие PL выполнено с $\mu$, равным этому инфимуму.
\end{examplebox}

\begin{trapbox}
Условие PL и выпуклость независимы. Функция из примера выше удовлетворяет PL, но не выпукла. Функция $f(x) = x^4$ выпукла, но условию PL не удовлетворяет: $\frac{1}{2}f'(x)^2 = 8x^6$ меньше $\mu x^4$ при любом $\mu > 0$, как только $|x|$ достаточно мало. Вблизи точки минимума функция $x^4$ слишком пологая для линейной скорости сходимости.
\end{trapbox}

\subsection{Число обусловленности и число итераций}\label{sub:rates-kappa}

При $\alpha = 1/L$ оценка имеет вид $f(x_k) - f^* \le (1 - 1/\kappa)^k\big(f(x_0) - f^*\big)$. Чтобы невязка уменьшилась до доли $\varepsilon$ от начальной, достаточно, чтобы $(1 - 1/\kappa)^k \le \varepsilon$, то есть
\handoutdisplay{
k \ge \frac{\ln \varepsilon}{\ln(1 - 1/\kappa)},
}
 где оба логарифма отрицательны. Так как $\ln(1 - 1/\kappa) \le -1/\kappa$, правая часть не больше $\kappa\ln(1/\varepsilon)$, а при больших $\kappa$ близка к этой величине. Число итераций растёт линейно по $\kappa$ и лишь логарифмически по точности: это и есть $O(\kappa\log(1/\varepsilon))$ из лекции.

\begin{examplebox}
Возьмём $\kappa = 1000$ и $\varepsilon = 10^{-6}$. Тогда $\ln\varepsilon/\ln(1 - 1/\kappa)$ после округления вверх даёт $k = 13809$, а оценка $\kappa\ln(1/\varepsilon)$ — примерно $13816$. Для квадратичной функции оптимальный шаг $2/(L + \mu)$ умножает каждую собственную компоненту ошибки на число, по модулю не большее $\rho = (\kappa - 1)/(\kappa + 1)$, поэтому за шаг невязка умножается не более чем на $\rho^2$, и для того же уменьшения нужно
\handoutdisplay{
k \ge \frac{\ln\varepsilon}{2\ln\rho},
}
 то есть $k = 3454$. Это примерно четверть числа итераций при шаге $1/L$, потому что $\ln\rho \approx -2/\kappa$ при больших $\kappa$.
\end{examplebox}

Эти числа — гарантии, и на деле метод может сойтись намного быстрее. В A2 при $\kappa = 2$ оценка даёт 20 итераций, а фактически нужно 10; в A3 при $\kappa = 50$ оценка даёт 684, а фактически нужно 245. Разницу объясняют два эффекта. Для квадратичной функции с шагом $\alpha = 1/L$ за каждый шаг $i$-я собственная компонента невязки умножается на $(1 - \lambda_i/L)^2$, так что самая медленная умножается на $(1 - \mu/L)^2$, то есть на квадрат гарантированного множителя; уже одно это вдвое сокращает число итераций. Кроме того, жёсткая компонента исчезает на первом же шаге, ведь её множитель равен $1 - L/L = 0$: в A3 имеем $f(x_0) = 25{,}5$, но $f(x_k) = \frac{1}{2}\,(0{,}98)^{2k}$ при $k \ge 1$.

С оптимальным шагом $2/(L + \mu) = \tfrac{2}{51}$ обе компоненты A3 за шаг умножаются на число, по модулю равное ровно $\tfrac{49}{51}$, оценка $f(x_k) \le \rho^{2k} f(x_0)$ выполняется как равенство, и оценка совпадает с фактическим числом итераций: 173 в обоих случаях.

\begin{selfcheckbox}
Что происходит с числом итераций при фиксированном $\kappa$, если заменить $\varepsilon$ на $\varepsilon^2$?
\selfcheckanswer Оно удваивается: число итераций пропорционально $\ln(1/\varepsilon)$, а $\ln(1/\varepsilon^2) = 2\ln(1/\varepsilon)$. Чтобы уменьшить невязку на вдвое большее число порядков, нужно вдвое больше итераций; при больших $\kappa$ вдвое большее число обусловленности тоже примерно удваивает число итераций.
\end{selfcheckbox}

% =====================================================================
\section{Автоматическое дифференцирование в обратном режиме}\label{sec:ad}

Градиентному спуску на каждом шаге нужен $\nabla f$. \textbf{Автоматическое дифференцирование} (AD) получает его из программы, которая вычисляет $f$: цепное правило (правило дифференцирования сложной функции) применяется к элементарным операциям по одной, и результат точен, если не считать ошибок округления. Обратный режим выдаёт весь градиент за время, в небольшое постоянное число раз превышающее стоимость одного вычисления $f$, сколько бы ни было переменных. Здесь всё делается вручную на небольших графах.

\subsection{Вычисление как граф}\label{sub:ad-graph}

Программа вычисляет $f$ через последовательность промежуточных переменных, каждая из которых — элементарная функция предыдущих. В примере из лекции $z = \sin(xy) + x^2$:
\handoutdisplay{
u = xy, \qquad v = \sin u, \qquad w = x^2, \qquad z = v + w.
}
 Вершины графа — входы $x, y$ и промежуточные переменные $u, v, w, z$; ребро $a \to b$ означает, что $b$ использует $a$. Каждая переменная использует только более ранние, поэтому циклов нет: это \textbf{ориентированный ациклический граф} (DAG). Порядок $x, y, u, v, w, z$, в котором программа вычисляет вершины, — \textbf{топологический порядок}: все рёбра в нём направлены вперёд.

Каждая вершина знает свои локальные частные производные: $\partial u/\partial x = y$, $\partial u/\partial y = x$, $\partial v/\partial u = \cos u$, $\partial w/\partial x = 2x$, $\partial z/\partial v = \partial z/\partial w = 1$. AD комбинирует эти числа по цепному правилу и никогда не строит формулу для $\partial z/\partial x$ целиком.

\handoutfigure{0.95}{fig_graph_c1}{Граф задачи C1 в точке $(x, y) = (2;\ -1)$. Значения на прямом проходе: $u = -2$, $v = -0{,}909297$, $w = 4$, $z = 3{,}090703$. Сопряжённые переменные обратного прохода: $\bar z = \bar v = \bar w = 1$, $\bar u = \cos u = -0{,}416147$, $\bar x = 4{,}416147$, $\bar y = -0{,}832294$.}

\subsection{Прямой режим}\label{sub:ad-forward}

Выберем направление $(\dot x, \dot y)$ в пространстве входов и будем вместе с каждым значением переносить его \textbf{касательное значение}: производную этого значения по выбранному направлению. Касательные значения распространяются по графу вперёд с теми же локальными производными:
\handoutdisplay{
\dot u = y\dot x + x\dot y, \qquad \dot v = \cos(u)\,\dot u, \qquad \dot w = 2x\,\dot x, \qquad \dot z = \dot v + \dot w.
}


Для отображения $F: \mathbb R^n \to \mathbb R^m$ с якобианом $J$ один прямой проход с начальным вектором $\dot x \in \mathbb R^n$ возвращает \textbf{произведение якобиана на вектор} $J\dot x$. Для скалярной $f$ это производная по направлению $\nabla f(x)^\top \dot x$, одно число, так что для полного градиента нужно $n$ проходов, по одному на каждый базисный вектор $e_1, \dots, e_n$.

\begin{examplebox}
В точке задачи C1 с начальным вектором $\dot x = 1$, $\dot y = 0$ один проход даёт $\dot z = y\cos u + 2x = 4{,}416147$ — частную производную по $x$. Для производной по $y$ нужен второй проход с $\dot x = 0$, $\dot y = 1$.
\end{examplebox}

\subsection{Обратный режим}\label{sub:ad-reverse}

Для скалярного выхода $z$ \textbf{сопряжённая переменная} вершины $v$ — это $\bar v = \partial z/\partial v$, производная выхода по этой вершине с учётом всего, что вычисляется из неё дальше. Начнём с $\bar z = 1$, все остальные сопряжённые переменные положим равными $0$, обойдём вершины в обратном топологическом порядке и для каждого ребра $u \to v$ прибавим вклад
\handoutdisplay{
\bar u \mathrel{+}= \bar v\, \frac{\partial v}{\partial u}.
}
 Когда обход доходит до вершины, все вершины, которые её используют, уже обработаны, поэтому её сопряжённая переменная уже вычислена полностью. Одной формулой:
\handoutdisplay{
\bar v_i = \sum_{j:\ i \to j} \bar v_j\, \frac{\partial v_j}{\partial v_i},
}
 где сумма берётся по всем вершинам $j$, использующим $v_i$.

В примере: $\bar z = 1$; вершина $z = v + w$ даёт $\bar v = 1$ и $\bar w = 1$; $w = x^2$ даёт $\bar x \mathrel{+}= 2x$; $v = \sin u$ даёт $\bar u = \cos u$; $u = xy$ даёт $\bar x \mathrel{+}= y\cos u$ и $\bar y \mathrel{+}= x\cos u$. Следовательно,
\handoutdisplay{
\bar x = y\cos u + 2x = \frac{\partial z}{\partial x}, \qquad \bar y = x\cos u = \frac{\partial z}{\partial y}.
}
 Один обратный проход дал обе частные производные.

Для $F: \mathbb R^n \to \mathbb R^m$ обратный проход, начатый с вектора $\bar z \in \mathbb R^m$, возвращает \textbf{произведение вектора на якобиан} $\bar z^\top J$. Для скалярной функции потерь ($m = 1$, $\bar z = 1$) это весь градиент за один проход, каким бы ни было $n$. Обратный режим, применённый к модели со скалярной функцией потерь и множеством параметров, — это \textbf{обратное распространение ошибки} (backpropagation).

\subsection{Почему сопряжённые переменные накапливаются}\label{sub:ad-accumulate}

Вершина, от которой зависят несколько других, получает вклад по каждому выходящему ребру, и цепное правило для функций многих переменных эти вклады складывает. Для $z = x^2 + x^3$ с вершинами $a = x^2$, $b = x^3$, $z = a + b$:
\handoutdisplay{
\frac{dz}{dx} = \frac{\partial z}{\partial a}\frac{da}{dx} + \frac{\partial z}{\partial b}\frac{db}{dx} = 2x + 3x^2.
}
 Поэтому правило и записано через $\mathrel{+}=$: сопряжённая переменная $\bar x$ собирает $2x$ от $a$ и $3x^2$ от $b$. Реализация, которая присваивает вместо того, чтобы прибавлять, сохраняет только вклад последнего обработанного ребра, то есть $2x$ или $3x^2$, и неверный ответ ещё и зависит от порядка обхода.

В общем случае $\partial z/\partial x$ — это сумма по всем путям из $x$ в $z$ произведений локальных производных вдоль пути. Обратный проход вычисляет эту сумму, не перечисляя путей: сопряжённая переменная каждой вершины уже содержит сумму по всем путям из этой вершины в $z$. В C1 вход $x$ используется в $u$ и $w$, поэтому у $\bar x$ два слагаемых; в C2 он входит и в $xy$, и в $xy + x$.

\begin{trapbox}
По той же причине фреймворки накапливают градиенты параметров от вызова к вызову. Перед вычислением нового градиента сохранённые сопряжённые переменные нужно обнулить, иначе новый градиент прибавится к старому.
\end{trapbox}

\subsection{Стоимость и память}\label{sub:ad-cost}

Пусть $f: \mathbb R^n \to \mathbb R$ вычисляется за $N$ элементарных операций. И прямой, и обратный проход выполняют ограниченный объём работы на каждую операцию, поэтому каждый из них в небольшое постоянное число раз дороже одного вычисления $f$. Режимы различаются числом проходов.

\begin{center}\small
\begin{tabular}{@{}>{\raggedright\arraybackslash\hspace{0pt}}p{0.17\linewidth}>{\raggedright\arraybackslash\hspace{0pt}}p{0.42\linewidth}>{\raggedright\arraybackslash\hspace{0pt}}p{0.29\linewidth}@{}}
\toprule
Свойство & Прямой режим & Обратный режим \\
\midrule
Начальный вектор & направление $\dot x \in \mathbb R^n$ & веса выходов $\bar z \in \mathbb R^m$ \\
Один проход даёт & $J\dot x$ — произведение якобиана на вектор & $\bar z^\top J$ — произведение вектора на якобиан \\
Градиент $f: \mathbb R^n \to \mathbb R$ & $n$ проходов & один прямой и один обратный проход \\
Якобиан $F: \mathbb R^n \to \mathbb R^m$ & $n$ проходов & $m$ проходов \\
Дополнительная память & по одному касательному значению на каждое используемое значение & все промежуточные значения (лента) \\
\bottomrule
\end{tabular}
\end{center}

Обратный режим дёшев по числу операций, но не по памяти. Локальным производным нужны значения прямого прохода ($\partial v/\partial u = \cos u$ требует $u$), поэтому прямой проход должен записать все промежуточные значения — \textbf{ленту}, — прежде чем начнётся обратный, и память растёт с длиной вычисления.

\textbf{Контрольные точки} (checkpointing) позволяют экономить память за счёт дополнительных вычислений. Прямой проход сохраняет только часть промежуточных значений — контрольные точки; во время обратного прохода каждый участок между двумя контрольными точками пересчитывается от предыдущей из них, когда понадобятся его значения. Если взять около $\sqrt N$ контрольных точек через каждые $\sqrt N$ операций, память сокращается до порядка $\sqrt N$ хранимых значений ценой примерно одного дополнительного прямого прохода.

\subsection{Обратное распространение для логистической функции потерь}\label{sub:ad-logistic}

Одно наблюдение состоит из признаков $x \in \mathbb R^d$ и метки $y \in \{0, 1\}$. Модель вычисляет
\handoutdisplay{
s = w^\top x + b, \qquad p = \sigma(s) = \frac{1}{1 + e^{-s}}, \qquad \ell = -y\log p - (1 - y)\log(1 - p).
}
 Так как $\log p = s - \log(1 + e^s)$ и $\log(1 - p) = -\log(1 + e^s)$, функция потерь равна $\ell = \log(1 + e^s) - ys$ — именно в таком виде она используется в C3.

Обратный проход начинается с $\bar\ell = 1$. Вершина $\ell(s)$ даёт
\handoutdisplay{
\bar s = \frac{d\ell}{ds} = \frac{e^s}{1 + e^s} - y = \sigma(s) - y,
}
 а вершина $s = w^\top x + b$ передаёт это значение параметрам: $\bar w = \bar s\, x$ и $\bar b = \bar s$. Итак,
\handoutdisplay{
\nabla_w \ell = \big(\sigma(s) - y\big)\,x, \qquad \frac{\partial\ell}{\partial b} = \sigma(s) - y,
}
 и значение $p = \sigma(s)$ с прямого прохода используется повторно. Для среднего значения потерь $\frac1n\sum_i \ell_i$ вклады наблюдений складываются и делятся на $n$: $\nabla_w \frac1n\sum_i \ell_i = \frac1n\sum_i \big(\sigma(s_i) - y_i\big)\,x_i$.

С метками $\tilde y = 2y - 1 \in \{-1, +1\}$ та же функция потерь записывается как
\handoutdisplay{
\ell = \log\big(1 + e^{-\tilde y s}\big).
}
 Проверим оба случая: при $y = 1$, $\tilde y = 1$ получаем $\log(1 + e^{-s}) = \log(1 + e^s) - s$, а при $y = 0$, $\tilde y = -1$ — $\log(1 + e^s)$. Производная $d\ell/ds = -\tilde y\,\sigma(-\tilde y s)$ в обоих случаях равна $\sigma(s) - y$. Два соглашения описывают одну и ту же функцию.

\begin{trapbox}
$\sigma(s) - y$ — производная по $s$, а не по $w$. В градиенте по $w$ есть множитель $x$; без него вместо вектора из $\mathbb R^d$ получается число.
\end{trapbox}

\subsection{Символьное, численное и автоматическое дифференцирование}\label{sub:ad-compare}

\begin{center}\small
\begin{tabular}{@{}>{\raggedright\arraybackslash\hspace{0pt}}p{0.15\linewidth}>{\raggedright\arraybackslash\hspace{0pt}}p{0.15\linewidth}>{\raggedright\arraybackslash\hspace{0pt}}p{0.24\linewidth}>{\raggedright\arraybackslash\hspace{0pt}}p{0.19\linewidth}>{\raggedright\arraybackslash\hspace{0pt}}p{0.13\linewidth}@{}}
\toprule
Метод & Основная идея & Точность & Типичная проблема & Стоимость градиента в $\mathbb R^n$ \\
\midrule
Символьное & преобразует формулу функции в формулу производной & точное & разбухание выражений: формула производной может быть намного больше программы & вычисление этой формулы, размер которой никак не контролируется \\
Конечные разности & сдвигает входы и делит разности значений на шаг & приближённое: в лучшем случае относительная ошибка около $\varepsilon^{2/3}$ для центральных разностей, где $\varepsilon$ — машинный эпсилон & при больших шагах — ошибка отбрасывания, при малых — потеря точности при вычитании близких чисел & $2n$ вычислений $f$ для центральных разностей \\
Автоматическое, обратный режим & применяет цепное правило к программе, вершина за вершиной & точное, если не считать ошибок округления & память под ленту; ветвления и циклы в программе, негладкие операции & в небольшое постоянное число раз дороже одного вычисления $f$ \\
\bottomrule
\end{tabular}
\end{center}

У центральной разности
\handoutdisplay{
\frac{\partial f}{\partial x_i}(x) \approx \frac{f(x + he_i) - f(x - he_i)}{2h}
}
 две ошибки. В формуле Тейлора $f(x \pm he_i) = f \pm h\,\partial_i f + \frac{h^2}{2}\partial_i^2 f \pm \frac{h^3}{6}\partial_i^3 f + O(h^4)$ при вычитании сокращаются чётные члены, поэтому ошибка \textbf{отбрасывания} равна $\frac{h^2}{6}|\partial_i^3 f| + O(h^3)$, то есть имеет порядок $h^2$. Каждое вычисленное значение $f$ содержит ошибку \textbf{округления} около $\varepsilon|f|$, где $\varepsilon$ — машинный эпсилон двойной точности. Разность двух близких чисел сохраняет эти абсолютные ошибки, и после деления на $2h$ получается ошибка порядка $\varepsilon|f|/h$. Суммарная ошибка
\handoutdisplay{
E(h) \approx \frac{h^2}{6}M + \frac{\varepsilon|f|}{h}, \qquad M = |\partial_i^3 f|,
}
 минимальна там, где $E'(h) = \frac{hM}{3} - \frac{\varepsilon|f|}{h^2} = 0$, то есть при $h^3 = 3\varepsilon|f|/M$. Поэтому лучший шаг имеет порядок $\varepsilon^{1/3} \approx 6 \cdot 10^{-6}$, а лучшая ошибка — порядок $\varepsilon^{2/3}$. У односторонней разности $\big(f(x + he_i) - f(x)\big)/h$ ошибка отбрасывания имеет порядок $h$, и то же рассуждение даёт лучший шаг порядка $\sqrt\varepsilon$.

Демонстрация D3 измеряет это для $\partial z/\partial x$ из C1 в точке $(x, y) = (2;\ -1)$. Относительная ошибка убывает как $h^2$ от $1{,}6 \cdot 10^{-4}$ при $h = 10^{-1}$, достигает минимума $5{,}7 \cdot 10^{-12}$ при $h = 10^{-5}$, близком к $\varepsilon^{1/3}$, и снова растёт до $1{,}4 \cdot 10^{-7}$ при $h = 10^{-10}$, когда начинает преобладать ошибка округления.

\handoutfigure{0.72}{fig_fd_error}{Относительная ошибка центральной разности для $\partial z/\partial x$ из C1 в точке $(x, y) = (2;\ -1)$ при $h$ от $10^{-1}$ до $10^{-10}$. При больших $h$ преобладает ошибка отбрасывания, и она убывает как $h^2$; при малых $h$ преобладает ошибка округления, и она растёт как $1/h$.}

У AD нет шага, который нужно выбирать, и нет ошибки отбрасывания: оно вычисляет точную производную программы в арифметике с плавающей точкой. Конечные разности остаются полезны как независимая проверка градиента: в нескольких точках, с $h$ около $\varepsilon^{1/3}$ и относительной мерой ошибки.

\subsection{Задачи блока C}\label{sub:ad-problems}

Три задачи семинара на обратный проход с полными решениями.

\problemhead{C1}{Граф из лекции в конкретной точке}{У доски}
\textbf{Условие.} Возьмите $z = \sin(xy) + x^2$ в точке $(x, y) = (2;\ -1)$ с промежуточными вершинами $u = xy$, $w = x^2$, $v = \sin u$. Выпишите вершины в топологическом порядке, выполните прямой проход, затем обратный проход, начиная с $\bar z = 1$. Найдите $z$, $\partial z/\partial x$ и $\partial z/\partial y$ и проверьте обе производные аналитически.

\textbf{Решение.}
\begingroup\sloppy
\begin{enumerate}[label=\textbf{\arabic*.}, leftmargin=*]
  \item Топологический порядок (каждая вершина идёт после своих входов): $x$, $y$, $u = xy$, $w = x^2$, $v = \sin u$, $z = v + w$. Из входа $x$ выходят два ребра: в $u$ и в $w$.
  \item Прямой проход в точке $(x, y) = (2;\ -1)$, все значения сохраняем:
\handoutdisplay{
u = -2,\quad w = 4,\quad v = \sin(-2) = -0{,}909297,\quad z = v + w = 3{,}090703.
}
 На обратном проходе снова понадобятся $u$, $x$ и $y$.
  \item Обратный проход, вершина $z = v + w$: начинаем с $\bar z = 1$. Обе локальные производные равны $1$, поэтому $\bar v = 1$ и $\bar w = 1$.
  \item Вершина $v = \sin u$, угол $u$ в радианах: $\bar u \mathrel{+}= \bar v \cos u = \cos(-2) = -0{,}416147$.
  \item Вершина $w = x^2$: локальная производная равна $2x$, поэтому $\bar x \mathrel{+}= \bar w \cdot 2x = 2 \cdot 2$. Это первый из двух вкладов в $\bar x$.
  \item Вершина $u = xy$: локальные производные равны $y$ и $x$, поэтому
\handoutdisplay{
\bar x \mathrel{+}= \bar u\,y = (-0{,}416147)(-1) = 0{,}416147, \qquad \bar y \mathrel{+}= \bar u\,x = (-0{,}416147) \cdot 2 = -0{,}832294.
}
 С учётом обоих вкладов $\bar x = 0{,}416147 + 2 \cdot 2 = 4{,}416147$; при перезаписи осталось бы только $0{,}416147$.
  \item Аналитическая проверка: $\partial z/\partial x = y\cos(xy) + 2x$ и $\partial z/\partial y = x\cos(xy)$ дают те же числа. Обратный проход вычислил эти формулы слагаемое за слагаемым: $\bar u\,y$ — это $y\cos(xy)$, а $\bar w \cdot 2x$ — это $2x$. Один проход дал обе частные производные.
\end{enumerate}
\endgroup
\begin{answerbox}
$z = 3{,}090703$, $\partial z/\partial x = 4{,}416147$, $\partial z/\partial y = -0{,}832294$.
\end{answerbox}

\problemhead{C2}{Входная переменная, использованная дважды}{Самостоятельно}
\textbf{Условие.} Возьмите $q = (xy + x)^2$ в точке $(x, y) = (2;\ 3)$ с вершинами $a = xy$, $b = a + x$, $q = b^2$. Вычислите $q$ прямым проходом, а $\partial q/\partial x$ и $\partial q/\partial y$ — одним обратным проходом, вручную, на бумаге.

\textbf{Решение.}
\begingroup\sloppy
\begin{enumerate}[label=\textbf{\arabic*.}, leftmargin=*]
  \item Топологический порядок: $x$, $y$, $a = xy$, $b = a + x$, $q = b^2$. Из $x$ выходят два ребра: одно в $a$, другое прямо в $b$.
  \item Прямой проход в точке $(x, y) = (2;\ 3)$: $a = xy = 6$, $b = a + x = 8$, $q = b^2 = 64$.
  \item Обратный проход: $\bar q = 1$, и вершина $q = b^2$ даёт $\bar b = \bar q \cdot 2b = 16$.
  \item Вершина $b = a + x$: обе локальные производные равны $1$, поэтому $\bar a \mathrel{+}= \bar b = 16$ и $\bar x \mathrel{+}= \bar b = 16$.
  \item Вершина $a = xy$: $\bar x \mathrel{+}= \bar a\,y = 48$ и $\bar y \mathrel{+}= \bar a\,x = 32$. Итого $\bar x = 16 + 48 = 64$; с \texttt{=} вместо \texttt{+=} получилось бы $48$ — только последнее слагаемое.
  \item Аналитическая проверка: $q = x^2(y + 1)^2$, поэтому $\partial q/\partial x = 2x(y + 1)^2 = 64$ и $\partial q/\partial y = 2x^2(y + 1) = 32$. Значение $q$ и производная $\partial q/\partial x$ совпадают здесь лишь потому, что $x = 2$: их отношение равно $x/2$.
\end{enumerate}
\endgroup
\begin{answerbox}
$q = 64$, $\partial q/\partial x = 48 + 16 = 64$, $\partial q/\partial y = 32$.
\end{answerbox}

\problemhead{C3}{Вершина логистической функции потерь}{Самостоятельно}
\textbf{Условие.} Для одного примера с логитом $s$ и меткой $y \in \{0, 1\}$ логистическая функция потерь равна $\ell(s) = \log(1 + e^s) - ys$, а $\sigma(s) = 1/(1 + e^{-s})$. При $s = 0{,}5$ и $y = 1$ разбейте $\ell$ на элементарные вершины, выполните прямой и обратный проходы и найдите $\ell$, $\sigma(s)$ и $d\ell/ds$. Затем перепишите функцию потерь для меток $\tilde y = 2y - 1 \in \{-1, 1\}$ как функцию одного только $\tilde y s$ и проверьте её для обеих меток.

\textbf{Решение.}
\begingroup\sloppy
\begin{enumerate}[label=\textbf{\arabic*.}, leftmargin=*]
  \item Вершины в топологическом порядке: $s$, $u = e^s$, $v = 1 + u$, $w = \log v$, $\ell = w - ys$. Метка $y$ — константа. Логит $s$ идёт в две вершины, $u$ и $\ell$, поэтому у $\bar s$ будет два слагаемых.
  \item Прямой проход при $s = 0{,}5$, $y = 1$:
\handoutdisplay{
\begin{gathered}
u = e^s = 1{,}648721,\quad v = 1 + 1{,}648721,\quad w = \log v = 0{,}974077, \\[6pt]
\ell = w - ys = 0{,}974077 - 1 \cdot 0{,}5 = 0{,}474077.
\end{gathered}
}

  \item Обратный проход: $\bar\ell = 1$. У вершины $\ell = w - ys$ локальная производная по $w$ равна $1$, а по $s$ равна $-y$, поэтому $\bar w = 1$ и $\bar s \mathrel{+}= -y = -1$.
  \item Через логарифм: $\bar v = \bar w / v = \dfrac{1}{1 + e^s}$, затем $\bar u = \bar v$, и наконец
\handoutdisplay{
\bar s \mathrel{+}= \bar u\,e^s = \frac{e^s}{1 + e^s} = \frac{1}{1 + e^{-s}} = \sigma(s) = 0{,}622459.
}

  \item Оба слагаемых вместе: $d\ell/ds = \bar s = \sigma(s) - y = 0{,}622459 - 1 = -0{,}377541$. Это $p - y$ из лекции, где $p = \sigma(s)$. Производная отрицательна, поэтому шаг градиентного спуска повышает логит положительного примера.
  \item Метки $\pm 1$: положим $\tilde y = 2y - 1$. При $y = 1$ получаем $\tilde y = 1$ и $\log(1 + e^s) - s = \log\dfrac{1 + e^s}{e^s} = \log(1 + e^{-s})$. При $y = 0$ получаем $\tilde y = -1$, а функция потерь остаётся прежней, $\log(1 + e^s)$. Оба случая записываются одной формулой:
\handoutdisplay{
\ell = \log\big(1 + e^{-\tilde y s}\big).
}

  \item Проверка в точке: $\tilde y = 1$ и $\log(1 + e^{-\tilde y s}) = 0{,}474077$ — то же число, что и $\ell$. Два соглашения задают одну и ту же функцию, поэтому у них совпадают все значения и производная $\sigma(s) - y$. В лекции функция потерь записана для меток из $\{0, 1\}$, а запись с $\pm 1$ — второе распространённое соглашение.
\end{enumerate}
\endgroup
\begin{answerbox}
$\ell = 0{,}474077$, $\sigma(s) = 0{,}622459$, $d\ell/ds = \sigma(s) - y = -0{,}377541$; при $\tilde y = 2y - 1$ функция потерь равна $\log(1 + e^{-\tilde y s})$.
\end{answerbox}

% =====================================================================
\section{Коллоквиум}\label{sec:colloquium}

На коллоквиуме сначала спрашивают определения, поэтому и здесь первым идёт список определений из первого раздела. За ним — четыре доказательства с ключевыми неравенствами и три задачи в формате билета с полными решениями.

\subsection{Список определений}\label{sub:colloquium-definitions}

Для каждого определения нужны название, формула и кванторы: для каких точек или пар точек формула должна выполняться.

\begin{itemize}
  \item Точка локального минимума: точка $x^*$, для которой найдётся $\varepsilon > 0$, при котором $f(x) \ge f(x^*)$ для всех $x$ с $\|x - x^*\| < \varepsilon$.
  \item Точка глобального минимума: точка $x^*$, для которой $f(x) \ge f(x^*)$ при всех $x$.
  \item Стационарная точка: точка $x$, в которой $f$ дифференцируема и $\nabla f(x) = 0$.
  \item Выпуклое множество: $(1 - t)x + ty \in C$ для всех $x, y \in C$ и $t \in [0, 1]$.
  \item Выпуклая функция (неравенство Йенсена): $f\big((1 - t)x + ty\big) \le (1 - t)f(x) + t f(y)$ для всех $x, y$ и $t \in [0, 1]$.
  \item Строго выпуклая функция: то же неравенство выполняется строго при всех $x \ne y$ и $t \in (0, 1)$.
  \item $\mu$-сильно выпуклая функция, $\mu > 0$: $f(y) \ge f(x) + \nabla f(x)^\top(y - x) + \frac\mu2\|y - x\|^2$ для всех $x, y$; или, что равносильно, функция $f - \frac\mu2\|\cdot\|^2$ выпукла.
  \item $L$-гладкая функция: $\|\nabla f(x) - \nabla f(y)\| \le L\|x - y\|$ для всех $x, y$.
  \item Число обусловленности: $\kappa = L/\mu \ge 1$; для $f = \frac{1}{2}x^\top Q x$ оно равно $\lambda_{\max}(Q)/\lambda_{\min}(Q)$.
  \item Направление спуска в точке $x$: вектор $d$, для которого $\nabla f(x)^\top d < 0$.
  \item Условие Поляка–Лоясевича (PL) с константой $\mu > 0$: $\frac{1}{2}\|\nabla f(x)\|^2 \ge \mu\big(f(x) - f^*\big)$ для всех $x$.
\end{itemize}

\subsection{Четыре доказательства}\label{sub:colloquium-proofs}

\begin{notebox}{Лемма о спуске}
Если $\nabla f$ липшицев с константой $L$, то $f(y) \le f(x) + \nabla f(x)^\top(y - x) + \frac L2\|y - x\|^2$. Положим $r(t) = f(x + t(y - x))$ и запишем $f(y) - f(x) = \int_0^1 r'(t)\,dt$, где $r'(t) = \nabla f(x + t(y - x))^\top(y - x)$. Вычтем $\nabla f(x)^\top(y - x)$ и оценим остаток с помощью неравенства Коши — Буняковского — Шварца и условия Липшица:
\handoutdisplay{
f(y) - f(x) - \nabla f(x)^\top(y - x) = \int_0^1 \big(\nabla f(x + t(y - x)) - \nabla f(x)\big)^\top (y - x)\,dt \le \int_0^1 Lt\|y - x\|^2\,dt = \frac L2\|y - x\|^2.
}

\end{notebox}

\begin{notebox}{Конечность бэктрекинга по Армихо}
Если $\nabla f$ липшицев с константой $L$ и $g = \nabla f(x) \ne 0$, то бэктрекинг вдоль $d = -g$ останавливается после конечного числа уменьшений шага. Лемма о спуске при $y = x - \alpha g$ даёт
\handoutdisplay{
f(x - \alpha g) \le f(x) - \alpha\Big(1 - \frac{L\alpha}{2}\Big)\|g\|^2,
}
 и правая часть не больше $f(x) - c\alpha\|g\|^2$, как только $1 - L\alpha/2 \ge c$, то есть при всех $\alpha \le 2(1 - c)/L$. Пробные шаги $\beta^j\alpha_0$ стремятся к $0$, поэтому один из них попадёт в этот промежуток. Принятый шаг равен либо $\alpha_0$, либо отвергнутому пробному шагу, умноженному на $\beta$, а любой отвергнутый пробный шаг больше $2(1 - c)/L$, поэтому $\alpha \ge \min\{\alpha_0,\ 2\beta(1 - c)/L\}$.
\end{notebox}

\begin{notebox}{Скорость сходимости для выпуклых функций}
Для выпуклой $L$-гладкой $f$ и $\alpha = 1/L$ выполнено $f(x_k) - f^* \le L\|x_0 - x^*\|^2/(2k)$. Неравенство для шага спуска $f(x_{k+1}) \le f(x_k) - \frac{1}{2L}\|g_k\|^2$, выпуклость $f(x_k) - f^* \le g_k^\top(x_k - x^*)$ и тождество трёх точек для $x_{k+1} = x_k - \frac1L g_k$ вместе дают ключевое неравенство
\handoutdisplay{
f(x_{k+1}) - f^* \le \frac L2\Big(\|x_k - x^*\|^2 - \|x_{k+1} - x^*\|^2\Big).
}
 Просуммируем его по первым $k$ шагам: правая часть телескопически сокращается. Остаётся воспользоваться тем, что $f(x_j)$ не возрастает.
\end{notebox}

\begin{notebox}{Линейная скорость при сильной выпуклости и при условии PL}
Из сильной выпуклости следует условие PL: минимизируем по $y$ нижнюю оценку $f(x) + \nabla f(x)^\top(y - x) + \frac\mu2\|y - x\|^2$ и получаем $f^* \ge f(x) - \frac{1}{2\mu}\|\nabla f(x)\|^2$. При условии PL и $\alpha = 1/L$ ключевое неравенство имеет вид
\handoutdisplay{
f(x_{k+1}) - f^* \le f(x_k) - f^* - \frac{1}{2L}\|\nabla f(x_k)\|^2 \le \Big(1 - \frac\mu L\Big)\big(f(x_k) - f^*\big),
}
 а применяя его шаг за шагом, получаем $(1 - \mu/L)^k$. Второе доказательство использует только условие PL, выпуклость в нём не нужна.
\end{notebox}

\subsection{Билет T1: обусловленность и скорость сходимости}\label{sub:colloquium-t1}

\begin{selfcheckbox}
Пусть $f(x) = \frac{1}{2}x^\top Q x$, где $Q = \operatorname{diag}(4;\ 100)$. Найдите $\mu$, $L$ и $\kappa$. При каких постоянных шагах градиентный спуск сходится из любой начальной точки? Сколько итераций нужно по гарантированной оценке при $\alpha = 1/L$, чтобы разность $f(x_k) - f^*$ уменьшилась до доли $10^{-6}$ от начального значения? Найдите оптимальный постоянный шаг, его коэффициент сжатия и число итераций, которое ему нужно для такого же уменьшения.
\selfcheckanswer $\mu = \lambda_{\min}(Q) = 4$, $L = \lambda_{\max}(Q) = 100$, $\kappa = L/\mu = 25$. В собственных координатах итерация имеет вид $x_{k+1,i} = (1 - \alpha\lambda_i)x_{k,i}$, поэтому она сходится из любой начальной точки ровно тогда, когда $|1 - \alpha\lambda_i| < 1$ для обоих $i$, то есть при $0 < \alpha < 2/L = 0{,}02$. При $\alpha = 1/L = 0{,}01$ в оценке $f(x_k) - f^* \le (1 - \mu/L)^k\big(f(x_0) - f^*\big)$ стоит множитель $1 - \mu/L = 0{,}96$, и неравенство $(0{,}96)^k \le 10^{-6}$ требует $k \ge \ln(10^{-6})/\ln(0{,}96)$, то есть $k = 339$. Оптимальный постоянный шаг минимизирует наибольший из множителей $\max\{|1 - \alpha\mu|,\ |1 - \alpha L|\}$. Первая величина убывает с ростом $\alpha$, а вторая растёт при $\alpha > 1/L$, поэтому минимум достигается там, где $1 - \alpha\mu = \alpha L - 1$:
\handoutdisplay{
\alpha^* = \frac{2}{L + \mu} = \tfrac{1}{52} \approx 0{,}019231, \qquad \rho = \frac{L - \mu}{L + \mu} = \frac{\kappa - 1}{\kappa + 1} = \tfrac{12}{13} \approx 0{,}923077.
}
 На каждом шаге обе компоненты $x_k$ по модулю умножаются ровно на $\rho$, поэтому $f(x_k) = \rho^{2k} f(x_0)$, и неравенство $\rho^{2k} \le 10^{-6}$ требует $k \ge \ln(10^{-6})/(2\ln\rho)$, то есть $k = 87$: почти вчетверо меньше итераций, чем по оценке при $1/L$.
\end{selfcheckbox}

\subsection{Билет T2: бэктрекинг по Армихо}\label{sub:colloquium-t2}

\begin{selfcheckbox}
Запустите бэктрекинг по Армихо для $f(x) = \frac{1}{2}x^\top \operatorname{diag}(2;\ 20)\,x$ из точки $x_0 = (1;\ 1)$ вдоль $d = -\nabla f(x_0)$ при $\alpha_0 = 1$, $\beta = 0{,}5$ и $c = 0{,}25$. Выпишите пробные шаги, найдите принятый шаг и число делений пополам, затем сверьте результат с теоремой о конечности бэктрекинга.
\selfcheckanswer Здесь $L = \lambda_{\max} = 20$, $f(x_0) = 11$, $g = \nabla f(x_0) = (2;\ 20)$ и $\|g\|^2 = 404$. Пробный шаг $\alpha$ принимается, если $f(x_0 - \alpha g) \le f(x_0) - c\alpha\|g\|^2$:
\handoutdisplay{
\begin{array}{c|c|c|c} \alpha & f(x_0 - \alpha g) & f(x_0) - c\alpha\|g\|^2 & \text{принят} \\ \hline 1 & 3611 & -90 & \text{нет} \\ 0{,}5 & 810 & -39{,}5 & \text{нет} \\ 0{,}25 & 160{,}25 & -14{,}25 & \text{нет} \\ 0{,}125 & 23{,}0625 & -1{,}625 & \text{нет} \\ 0{,}0625 & 1{,}390625 & 4{,}6875 & \text{да} \end{array}
}
 Цикл останавливается после 4 делений пополам; принят шаг $\alpha = 0{,}0625$. Теорема гарантирует, что принимается любой шаг $\alpha \le 2(1 - c)/L = 0{,}075$, поэтому цикл не пойдёт дальше первого пробного шага, меньшего этого значения, а это $0{,}0625$. Поскольку хотя бы один пробный шаг был отвергнут, принятый шаг удовлетворяет и оценке $\alpha \ge \beta \cdot 2(1 - c)/L = 0{,}0375$. Обе проверки согласуются с расчётом: $0{,}0375 \le 0{,}0625 \le 0{,}075$.
\end{selfcheckbox}

\subsection{Билет T3: обратный проход}\label{sub:colloquium-t3}

\begin{selfcheckbox}
Вычислите $f(x, y) = x e^y + y^2$ и обе частные производные в точке $(x, y) = (2;\ 1)$ обратным проходом: выпишите вершины в топологическом порядке, найдите значения на прямом проходе, затем сопряжённые переменные и сверьте результат с аналитическим градиентом.
\selfcheckanswer Вершины в топологическом порядке: входы $x, y$, затем $a = e^y$, $b = xa$, $c = y^2$, $f = b + c$. Значения на прямом проходе: $a = e \approx 2{,}718282$, $b = 5{,}436564$, $c = 1$, $f = 6{,}436564$. Обратный проход начинается с $\bar f = 1$: вершина $f = b + c$ даёт $\bar b = 1$ и $\bar c = 1$; $c = y^2$ даёт $\bar y \mathrel{+}= \bar c \cdot 2y = 2$; $b = xa$ даёт $\bar x \mathrel{+}= \bar b\, a = 2{,}718282$ и $\bar a \mathrel{+}= \bar b\, x = 2$; $a = e^y$ даёт $\bar y \mathrel{+}= \bar a\, e^y = 5{,}436564$. Вход $y$ используется в двух вершинах, поэтому его сопряжённая переменная накапливается: $\bar y = 2 + 5{,}436564 = 7{,}436564$. Аналитически $\partial f/\partial x = e^y = e$ и $\partial f/\partial y = x e^y + 2y = 2e + 2$; в этой точке они равны $2{,}718282$ и $7{,}436564$. Если перезаписывать $\bar y$, а не накапливать, получится $2$ или $5{,}436564$ в зависимости от порядка.
\end{selfcheckbox}

% =====================================================================
\section{Полезные константы}\label{sec:constants}

Для шага $1/L$, границы устойчивости $2/L$ и скорости $1 - \mu/L$ нужны константы задачи. У квадратичной функции это собственные значения, а для логистической регрессии у $L$ есть простая верхняя оценка.

\subsection{Логистическая регрессия}\label{sub:constants-logistic}

Данные — векторы $x_1, \dots, x_n \in \mathbb R^d$, строки матрицы $X \in \mathbb R^{n \times d}$, с метками $\tilde y_i \in \{-1, +1\}$. Средняя логистическая функция потерь равна
\handoutdisplay{
f(w) = \frac1n\sum_{i=1}^n \log\big(1 + e^{-\tilde y_i w^\top x_i}\big),
}
 или, в записи с метками $\{0, 1\}$, $\frac1n\sum_i\big(\log(1 + e^{s_i}) - y_i s_i\big)$, где $s_i = w^\top x_i$ и $\tilde y_i = 2y_i - 1$; это одна и та же функция. Свободный член включается в $w$, если добавить к каждому $x_i$ постоянный признак $1$.

\begin{theorembox}{Утверждение: Гладкость логистической функции потерь.}
Обозначим $\sigma_i = \sigma(w^\top x_i)$. Тогда
\handoutdisplay{
\nabla^2 f(w) = \frac1n X^\top D X, \qquad D = \operatorname{diag}\big(\sigma_i(1 - \sigma_i)\big),
}
 и $0 \preceq \nabla^2 f(w) \preceq \frac{1}{4n}X^\top X$ при всех $w$. Значит, функция $f$ выпуклая и $L$-гладкая, причём
\handoutdisplay{
L \le \frac{\lambda_{\max}(X^\top X)}{4n}.
}

\end{theorembox}

\begin{proofbox}
Если метки взяты из $\{0, 1\}$, то $i$-е слагаемое равно $\varphi(s_i) - y_i s_i$, где $\varphi(s) = \log(1 + e^s)$, $\varphi'(s) = \sigma(s)$ и $\varphi''(s) = \sigma(s)\big(1 - \sigma(s)\big)$; линейное слагаемое кривизны не даёт. По цепному правилу с $s_i = w^\top x_i$ гессиан $i$-го слагаемого равен $\sigma_i(1 - \sigma_i)\,x_i x_i^\top$, а среднее таких матриц — $\frac1n X^\top D X$. Поскольку $p(1 - p) \le \frac{1}{4}$ при $p \in [0, 1]$, причём равенство достигается при $p = \frac{1}{2}$, для любого $v$
\handoutdisplay{
0 \le v^\top \nabla^2 f(w)\, v = \frac1n\sum_i \sigma_i(1 - \sigma_i)(x_i^\top v)^2 \le \frac{1}{4n}\sum_i (x_i^\top v)^2 = \frac{1}{4n}\|Xv\|^2 \le \frac{\lambda_{\max}(X^\top X)}{4n}\|v\|^2.
}
 Дважды дифференцируемая функция, у которой всюду $0 \preceq \nabla^2 f \preceq L I$, выпукла, а её градиент липшицев с константой $L$, потому что $\nabla f(y) - \nabla f(x) = \int_0^1 \nabla^2 f\big(x + t(y - x)\big)(y - x)\,dt$.
\end{proofbox}

Здесь $\lambda_{\max}(X^\top X) = \|X\|_2^2$ — квадрат наибольшего сингулярного числа $X$. Оценка достигается в точке $w = 0$, где все $\sigma_i = \frac{1}{2}$ и $\nabla^2 f(0) = \frac{1}{4n}X^\top X$, поэтому меньшей константы, верной на всём пространстве, не существует. Функция потерь не является сильно выпуклой: вдоль луча $w = tv$, где $x_i^\top v \ne 0$ при всех $i$, веса $\sigma_i(1 - \sigma_i)$ стремятся к $0$ при $t \to \infty$, а вместе с ними и гессиан.

Добавим штраф: $f_\lambda(w) = f(w) + \lambda\|w\|^2$ с весом $\lambda > 0$ (здесь это число, а не собственное значение). Гессиан получает добавку $\nabla^2\big(\lambda\|w\|^2\big) = 2\lambda I$:
\handoutdisplay{
2\lambda I \preceq \nabla^2 f_\lambda(w) \preceq \Big(\frac{\lambda_{\max}(X^\top X)}{4n} + 2\lambda\Big) I.
}
 Значит, $f_\lambda$ — $2\lambda$-сильно выпуклая и $L$-гладкая функция, причём $L \le \frac{\lambda_{\max}(X^\top X)}{4n} + 2\lambda$, и
\handoutdisplay{
\kappa = \frac L\mu \le \frac{\lambda_{\max}(X^\top X)/(4n) + 2\lambda}{2\lambda} = 1 + \frac{\lambda_{\max}(X^\top X)}{8n\lambda}.
}
 Тогда градиентный спуск с $\alpha = 1/L$ сходится линейно, а чем меньше $\lambda$, тем больше оценка для $\kappa$ и тем медленнее гарантированная скорость.

\begin{trapbox}
В некоторых источниках штраф записывают как $\frac\lambda2\|w\|^2$. Тогда гессиан получает добавку $\lambda I$, а не $2\lambda I$, задача становится $\lambda$-сильно выпуклой, и соответственно меняются все константы выше.
\end{trapbox}

\subsection{Квадратичные функции семинара}\label{sub:constants-quadratics}

Для $f(x) = \frac{1}{2}x^\top Q x$ с диагональной матрицей $Q$: $\mu = \lambda_{\min}(Q)$, $L = \lambda_{\max}(Q)$, градиентный спуск сходится при $0 < \alpha < 2/L$, линейная оценка скорости верна при $\alpha = 1/L$, а оптимальный постоянный шаг равен $2/(L + \mu)$.

\begin{center}\small
\begin{tabular}{@{}llrrrrrr@{}}
\toprule
Задача & $Q$ & $\mu$ & $L$ & $\kappa$ & $1/L$ & $2/L$ & $2/(L + \mu)$ \\
\midrule
A1 & $\operatorname{diag}(1;\ 10)$ & $1$ & $10$ & $10$ & $0{,}1$ & $0{,}2$ & $\tfrac{2}{11}$ \\
A2 & $\operatorname{diag}(1;\ 2)$ & $1$ & $2$ & $2$ & $0{,}5$ & $1$ & $\tfrac{2}{3}$ \\
A3 & $\operatorname{diag}(1;\ 50)$ & $1$ & $50$ & $50$ & $0{,}02$ & $0{,}04$ & $\tfrac{2}{51}$ \\
\bottomrule
\end{tabular}
\end{center}

В задаче B1 используется матрица из A1, в B2 и B3 — матрица из A3.

\end{document}
